% Build with: tectonic --keep-intermediates main.tex
\documentclass[11pt]{article}
\usepackage[margin=27mm]{geometry}
\usepackage[T1]{fontenc}
\usepackage{lmodern}
\usepackage{amsmath,amssymb,amsthm,mathtools}
\usepackage{booktabs,array,microtype,float,etoolbox}
\usepackage[shortlabels]{enumitem}
\usepackage[numbers,sort&compress]{natbib}
\usepackage[colorlinks=true,allcolors=blue,hypertexnames=false]{hyperref}
\usepackage{upquote}
\usepackage{doi}

\newtheorem{mainthm}{Theorem}
\renewcommand{\themainthm}{\Alph{mainthm}}
\newtheorem{theorem}{Theorem}[section]
\newtheorem{proposition}[theorem]{Proposition}
\newtheorem{lemma}[theorem]{Lemma}
\newtheorem{corollary}[theorem]{Corollary}
\theoremstyle{definition}
\newtheorem{definition}[theorem]{Definition}
\newtheorem{example}[theorem]{Example}
\theoremstyle{remark}
\newtheorem{remark}[theorem]{Remark}
\numberwithin{equation}{section}
% a statement is not broken across pages
\AtBeginEnvironment{mainthm}{\samepage}
\AtBeginEnvironment{theorem}{\samepage}
\AtBeginEnvironment{proposition}{\samepage}
\AtBeginEnvironment{lemma}{\samepage}
\AtBeginEnvironment{corollary}{\samepage}
\AtBeginEnvironment{definition}{\samepage}

\newcommand{\CC}{\mathbb C}
\newcommand{\ZZ}{\mathbb Z}
\newcommand{\Wclass}{\mathcal W}
\newcommand{\hatt}[1]{\widehat{#1}}
\newcommand{\cyl}{\overline{\CC/\Gamma}}
\newcommand{\Bil}{\mathcal B}
\DeclareMathOperator{\Res}{Res}
\DeclareMathOperator{\sech}{sech}
\DeclareMathOperator{\csch}{csch}
\DeclareMathOperator{\cs}{cs}
\DeclareMathOperator{\ds}{ds}
\DeclareMathOperator{\ns}{ns}
\DeclareMathOperator{\sn}{sn}
\DeclareMathOperator{\cn}{cn}
\DeclareMathOperator{\dn}{dn}
\DeclareMathOperator{\dist}{dist}
\setlength{\emergencystretch}{2em}
\setlist{itemsep=2pt,topsep=4pt}

\title{Meromorphic travelling waves for pure-power equations:\\
one pole orbit and a two-wave bound}
\author{Alexander Migita}
\date{1 October 2026}

\begin{document}
\maketitle


\begin{abstract}
The equation \(P(d/dz)\,v+v^n/n=0\), where \(P\) is a polynomial of
degree \(p\) with complex coefficients and \(n\ge2\), describes the
travelling waves of evolution equations with a pure-power nonlinearity
and arbitrary linear dispersion and dissipation. We prove that, when
\(p\ge3\), every nonconstant solution meromorphic in the whole complex
plane has exactly one pole modulo its twisted periods, the translations
that multiply it by a root of unity. Hence it is rational, rational in
an exponential, or elliptic, and lies in an explicit
finite-dimensional space. No condition on the coefficients is imposed;
in particular, the equation need not determine the Laurent series at a
pole. Up to translation and multiplication by a root of unity, one
equation has at most two such solutions, and two do occur. It has at
most one unless \(n\le3\) and \(P\) is even. Once the twisted
periods of a solution are known, its power \(v^{n-1}\) is determined by
\(P\); for elliptic solutions this gives an explicit necessary
condition on the coefficients.
\end{abstract}

\noindent\textbf{Keywords.} Meromorphic travelling waves; exact
solutions; twisted periods; Laurent series; elliptic functions of the
second kind.

\noindent\textbf{MSC2020.} 34M05, 35C07, 33E05.

% ======================================================================
\section{Introduction}\label{sec:intro}

\subsection{The equation and its waves}

Exact travelling waves of nonlinear evolution equations are usually
found with an ansatz: a polynomial in \(\tanh\), a combination of
\(\sech\) and \(\tanh\), a Jacobi or Weierstrass elliptic function, or a
rational function of an exponential
\citep{MusetteConte2003,Kudryashov2005Simplest,Kudryashov2010}. Every
such ansatz produces a function that is meromorphic in the whole
complex plane and belongs to the class
\[
 \Wclass=\{\text{rational functions of } z\}
 \cup\{\text{rational functions of } e^{kz},\ k\ne0\}
 \cup\{\text{elliptic functions}\}.
\]
Two questions arise. Can a travelling wave be meromorphic in \(\CC\)
without belonging to \(\Wclass\)? And which functions of \(\Wclass\),
and how many of them, solve one given equation?

We study both questions for the ordinary differential equation
\begin{equation}\label{eq:main}
 P(D)v+\frac{v^n}{n}=0,\qquad D=\frac{d}{dz},
\end{equation}
where \(n\ge2\) and \(P\) is a monic polynomial of degree \(p\) with
complex coefficients. For \(p\ge3\) the answer to the first question is
no: all nonconstant solutions that are meromorphic in \(\CC\) belong to
\(\Wclass\). For odd \(p\) this follows from a theorem of
\citet[Theorem~2]{Eremenko2006}; for even \(p\) it was known for
\(P(D)=D^p+a_0\) \citep{EremenkoLiaoNg2009} and for the even operators
of degree four \citep{DeminaKudryashovKawahara2010,ConteNgWong2012}.
As to the second, these solutions lie in explicit
finite-dimensional spaces, and up to the symmetries of the equation
there are at most two of them.

Equation~\eqref{eq:main} is the travelling-wave equation of evolution
equations such as
\begin{equation}\label{eq:pde}
 U_t+U^{n-1}U_x+\sum_{j=1}^{p}\alpha_j\,\partial_x^{\,j+1}U=0,
 \qquad \alpha_p\ne0.
\end{equation}
Indeed, travelling waves \(U(x,t)=w(x-st)\) of \eqref{eq:pde} satisfy,
after one integration,
\[
 \sum_{j=1}^{p}\alpha_jw^{(j)}-sw+\frac{w^n}{n}+C=0
\]
with a constant \(C\). When \(C=0\), the substitution
\(w=\alpha_p^{1/(n-1)}v\) gives \eqref{eq:main} with
\(P(\rho)=\alpha_p^{-1}\bigl(\sum_j\alpha_j\rho^j-s\bigr)\). So the
speed \(s\) enters the constant term of \(P\), and \(P\) is an even
polynomial exactly when \eqref{eq:pde} contains only derivatives of odd
order, that is, when the evolution equation has no dissipative terms.
For \(n=2\) a shift of \(w\) removes the constant \(C\) and changes
only the constant term of \(P\). For \(n\ge3\) we take \(C=0\); when
\(C\ne0\), moving a constant solution to \(0\) produces lower powers of
\(v\), and Appendix~\ref{app:general} treats such nonlinearities.

For example, the Kuramoto--Sivashinsky
\citep{KuramotoTsuzuki1976,Sivashinsky1977} and Kawahara
\citep{Kawahara1972} equations and the Nikolaevskiy equation with
dispersion \citep{KudryashovMigita2007} give \(n=2\) with \(p=3\),
\(p=4\) and \(p=5\) respectively, and the Swift--Hohenberg equation
\citep{SwiftHohenberg1977} gives \(n=3\) with \(p=4\). The last
equation, like reaction--diffusion equations, is not of the form
\eqref{eq:pde}: its travelling-wave equation is used as it stands,
without an integration, so \(p\) equals the order of the equation, and the speed enters as the
coefficient of \(D\) (Section~\ref{sec:sh}).

Throughout the paper, a \emph{wave} is a nonconstant solution of
\eqref{eq:main} that is meromorphic in the whole complex plane.
Suppose a wave \(v\) has a pole of order \(q\) at \(z_0\), with leading
coefficient \(c\). The highest derivative \(v^{(p)}\) must balance the
nonlinear term \(v^n/n\), which gives
\begin{equation}\label{eq:balance}
 p=(n-1)q,\qquad c^{\,n-1}=-n(-1)^p(q)_p,
\end{equation}
where \((a)_k=a(a+1)\cdots(a+k-1)\) is the rising factorial. A solution
without poles is constant (Proposition~\ref{prop:existence}). Hence
waves exist only when \(n-1\) divides \(p\). All poles of all waves
then have the same order \(q=p/(n-1)\), and the \(n-1\) possible leading
coefficients differ by factors \(\omega\) with \(\omega^{n-1}=1\).

Equation~\eqref{eq:main} is invariant under translations and under the
\emph{phase changes} \(v\mapsto\omega v\), where \(\omega\) lies in the
group \(\mu_{n-1}\) of \((n-1)\)st roots of unity. Combining the two
gives the notion around which this paper is organized.

\begin{definition}\label{def:twisted}
A \emph{twisted period} of a wave \(v\) is a number \(\gamma\in\CC\)
such that
\[
 v(z+\gamma)=\chi(\gamma)\,v(z)\qquad\text{for some }
 \chi(\gamma)\in\mu_{n-1}.
\]
The number \(\chi(\gamma)\) is the \emph{multiplier} of \(v\) under the
translation \(z\mapsto z+\gamma\). The twisted periods form a group
\(\Gamma\), and \(\chi\colon\Gamma\to\mu_{n-1}\) is a homomorphism, the
\emph{character} of \(v\). Its kernel is the group of ordinary periods.
\end{definition}

For instance, \(\sech z\) changes sign under \(z\mapsto z+i\pi\),
whereas \(\tanh z\) is unchanged: \(i\pi\) is a twisted period of
\(\sech\) with multiplier \(-1\), and an ordinary period of \(\tanh\).
Both functions have a single pole modulo \(i\pi\); modulo its ordinary
period \(2\pi i\), \(\sech z\) has two poles, with opposite residues.
Twisted periods that are not ordinary periods can occur only for
\(n\ge3\), since \(\mu_1=\{1\}\).

\subsection{Results}

In the statements we write \(m=n-1\) and \(q=p/m\).

\begin{mainthm}[One pole orbit]\label{thm:A}
Let \(n\ge2\) and let \(P\) have degree \(p\ge3\). Every wave \(v\)
of \eqref{eq:main} has exactly one pole modulo its twisted periods: if
\(z_0\) is a pole, then the set of all poles is the orbit
\(z_0+\Gamma\). After the translation that moves \(z_0\) to \(0\), the
wave has one of the following forms.
\begin{enumerate}[label=\textup{(\roman*)},leftmargin=2.2em]
\item \(\Gamma=\{0\}\), and \(v\) is a polynomial of degree \(q\) in
\(1/z\) without constant term.
\item \(\Gamma=(2\pi i/k)\ZZ\) for some \(k\ne0\). If
\(e^{2\pi ir/m}\), with \(1\le r\le m\), is the multiplier of \(v\)
under \(z\mapsto z+2\pi i/k\), then
\begin{equation}\label{eq:periodic-form}
 v=\beta+\sum_{j=0}^{q-1}d_j\,D^jg_r,\qquad
 g_r(z)=\frac{k\,e^{rkz/m}}{e^{kz}-1},
\end{equation}
with constants \(\beta,d_0,\dots,d_{q-1}\), where \(d_{q-1}\ne0\), and
\(\beta=0\) unless \(r=m\).
\item \(\Gamma\) is a lattice, and \(v\) is an elliptic function of the
second kind: \(v(z+\gamma)=\chi(\gamma)v(z)\) for \(\gamma\in\Gamma\),
with poles of order \(q\) at the points of \(\Gamma\) and nowhere else.
\end{enumerate}
In particular, every wave belongs to \(\Wclass\).
\end{mainthm}

The last assertion holds because the ordinary periods form a subgroup
of index at most \(m\) in \(\Gamma\): the character maps \(\Gamma\)
into \(\mu_m\). So a wave has at most \(m=n-1\) poles modulo its
ordinary periods; in case~(ii) it is a rational function of
\(e^{kz/m}\), and in case~(iii) it is elliptic for a sublattice of
\(\Gamma\). By the forms (i)--(iii), the search for all waves of a
given equation reduces to substitutions with finitely many unknown
coefficients; Section~\ref{sec:sh} carries one out.

In \eqref{eq:periodic-form}, the function \(g_r\) has simple poles
exactly at the points of \((2\pi i/k)\ZZ\), with residue \(1\) at
\(0\); it is multiplied by \(e^{2\pi ir/m}\) under
\(z\mapsto z+2\pi i/k\), and it stays bounded as
\(\operatorname{Re}(kz)\to\pm\infty\). Its logarithmic derivative is
\(rk/m-G\), where
\[
 G=g_m=\frac{k\,e^{kz}}{e^{kz}-1}=\frac k2\Bigl(1+\coth\frac{kz}{2}\Bigr),
 \qquad DG=G(k-G).
\]
By induction, \(D^jg_r\) is \(g_r\) times a polynomial of degree \(j\)
in \(G\); so \eqref{eq:periodic-form} says that \(v=\beta+g_r\,Q(G)\)
for a polynomial \(Q\) of degree \(q-1\). For \(n=2\) this is a
polynomial of degree \(q\) in \(G\), and the translation
\(z\mapsto z+i\pi/k\) turns it into a polynomial in \(\tanh(kz/2)\),
the form assumed by the tanh method. For \(n=3\) and \(r=1\), the same
translation turns \(g_1\) into a multiple of \(\sech(kz/2)\), and \(v\)
into \(\sech\) times a polynomial in \(\tanh\). In case~(iii) the space
of such elliptic functions has an explicit basis
(Lemma~\ref{lem:second-kind}).

In case~(ii) we call \(\operatorname{Re}(kz)\to\pm\infty\) the two
\emph{ends} of the wave (for real \(k\) these are \(z\to\pm\infty\) along
the real axis) and \(k\) its \emph{rate}; replacing \(k\) by \(-k\)
interchanges the two ends. A wave is \emph{twisted} if its character
is nontrivial, that is, if some twisted period is not an ordinary
period; otherwise it is \emph{untwisted}.

\begin{mainthm}[At most two waves]\label{thm:B}
Let \(n\ge2\) and let \(P\) have degree \(p\ge3\). Up to translation
and phase change, equation~\eqref{eq:main} has at most two waves. It
has at most one unless \(n\le3\) and \(P\) is even. For
\(P(D)=D^4-s\) with \(s\ne0\) it has exactly two, both for \(n=2\) and
for \(n=3\).
\end{mainthm}

The theorem counts all waves, complex and unbounded ones included. It
bounds their number; which operators admit a wave is a different
question (Section~\ref{sec:scope}). In terms of \eqref{eq:pde}: at a
fixed speed and a fixed integration constant, an evolution equation of
order \(p+1\ge4\) has at most two meromorphic travelling waves up to
translation and phase change, and at most one if it contains a
dissipative term or if \(n\ge4\). Here the constant is zero when
\(n\ge3\). For \(n=2\) it is arbitrary, but it has to be fixed,
because \(P\) depends on it (Section~\ref{sec:ks}).
Section~\ref{sec:example} exhibits the two waves of \(D^4-s\) for
\(n=3\); one of them is twisted.

The third result shows how much of a wave the operator determines
before the wave is known. We normalize a wave so that it has a pole at
\(0\) with leading coefficient \(c\), where \(c\) is a fixed root of
\eqref{eq:balance}, and put \(u=v/c=z^{-q}+O(z^{1-q})\).

\begin{mainthm}[The power of a wave]\label{thm:C}
Let \(p\ge3\). The coefficients of \(z^{-p},\dots,z^{-1},z^{0}\) in
the Laurent series of \(u^{\,n-1}\) at \(0\) are polynomials
\(\ell_0=1,\ell_1,\dots,\ell_p\) in the coefficients of \(P\); they
are the same for all waves. Together with the group \(\Gamma\) they
determine \(u^{\,n-1}\) by the explicit formula of
Theorem~\ref{thm:power}. In particular, the residue
\(\mathcal R(P)=\ell_{p-1}\) of \(u^{\,n-1}\) at \(0\) satisfies
\[
 \mathcal R(P)=0\ \text{ if \(\Gamma\) is a lattice},\qquad
 u_+^{\,n-1}-u_-^{\,n-1}=k\,\mathcal R(P)\ \text{ if }
 \Gamma=(2\pi i/k)\ZZ,
\]
where \(u_\pm\) are the limits of \(u\) at the ends
\(\operatorname{Re}(kz)\to\pm\infty\).
\end{mainthm}

For example, the Kuramoto--Sivashinsky equation with dispersion,
\(U_t+UU_x+\mu U_{xx}+\delta U_{xxx}+\nu U_{xxxx}=0\), has
\(\mathcal R=(16\mu\nu-\delta^2)/(608\nu^2)\); so its elliptic
travelling waves require \(\delta^2=16\mu\nu\)
(Section~\ref{sec:ks}).

Three further results follow from the theorems
(Sections~\ref{sec:power} and~\ref{sec:consequences}). If \(P\) is
even, every wave is even or odd about each of its poles, and a
\emph{front}, that is, a wave with different limits at its two ends,
exists only for odd \(q\). If \(n\ge3\) and \(P\) is an even polynomial of degree
\(p\ge4\) plus a nonzero multiple of \(D\), there are no waves at all.
And two waves of one equation that are not related by translation and
phase change never have the same twisted periods.

The theorems need \(p\ge3\). For \(p=2\) one equation can have
infinitely many waves, and some equations have waves outside
\(\Wclass\). The one-pole conclusion still holds for \(p=1\),
for even \(P\) of degree two, and for all elliptic waves with \(p=2\)
(Appendix~\ref{sec:low-order}).

\paragraph{The proofs.} The proofs use Laurent series, Cauchy's
estimate, and the identity theorem. At a pole, the Laurent series of a
wave is determined by \(P\) up to a root of unity and, when \(p\) is
even, one number \(h\), the \emph{free coefficient}. If the series at
two poles agree up to a root of unity, the difference of the poles is a
twisted period. So a wave whose poles all carry the same \(h\) has one
pole modulo \(\Gamma\), and two waves that share a value of \(h\)
coincide after a translation and a phase change. To find the values of
\(h\) we form combinations of a wave and its derivatives that have no
poles. By Cauchy's estimate a wave is large only near its poles; so
such a combination is constant as soon as it takes finitely many
values at the poles. Its Taylor coefficients of positive order at a
pole then vanish, and they are polynomials in \(h\) whose coefficients
depend only on \(P\). This gives an equation for \(h\) of degree one
when \(n\ge4\) and of degree two when \(n\le3\), and a further equation
of degree one when \(P\) is not even. A count of zeros and poles shows
that one wave cannot carry both roots of the quadratic equation.
Section~\ref{sec:example} carries out the argument on one equation, and
Appendix~\ref{app:notation} lists the notation.

\subsection{Earlier work}\label{sec:related}

\emph{Odd degree.} For odd \(p\) there is no free coefficient, and
\citet[Theorem~2]{Eremenko2006} showed by Nevanlinna theory
\citep{Laine1993,ConteNgWu2018} that all meromorphic solutions belong
to \(\Wclass\), with at most \(n-1\) poles per period; the argument
that two poles with equal Laurent series differ by a period is his
\citep{Eremenko1982,Eremenko2006}. For \(n=2\) the forms (i)--(iii) are
those of \citet[Theorem~1]{DeminaKudryashov2011}, who also note that an
equation with one Laurent series has at most one meromorphic solution
up to translation. The case \(n=4\), \(p=3\), with lower powers in the
nonlinearity, is due to \citet{ConteNg2010}.

\emph{Even degree.} The free coefficient was noted by
\citet{Eremenko1982}; Lemma~\ref{lem:resonance} is Proposition~2.1 of
\citet{Wu2014}, which covers every equation with the dominant terms
\(v^{(p)}\) and \(v^n\). For \(P(D)=D^p+a_0\) (in their setting, \(D^p\) with
a general nonlinearity), \citet*[Lemma~4]{EremenkoLiaoNg2009} fixed
this coefficient by the first integral that comes from the integrating
factor \(v'\). In the same way all meromorphic solutions were found
for even operators of degree four: for \(n=2,3,5\) by
\citet{DeminaKudryashovKawahara2010}, for \(n=3\) with a general cubic
nonlinearity by \citet*[Theorem~1]{ConteNgWong2012}, and for \(n=5\)
with a cubic term by \citet[Theorem~2.1]{Wu2014}. Their lists contain
the waves that we call twisted, and the two-wave examples of
Sections~\ref{sec:example} and~\ref{sec:n2}. The same argument works
for every even operator (Appendix~\ref{app:general}); for an operator
that is not even this integral is not available. For \(n=2\) and
operators of degree four that are not even, \citet{Demina2022}
classified the solutions that are elliptic or rational in an
exponential, and each of them has one pole per period
\citep[Section~5.6, Theorem~5.8]{BelovaThesis}.

\emph{Examples and tools.} The Swift--Hohenberg waves of
Section~\ref{sec:sh} were found by
\citet{KudryashovSinelshchikov2012SH}, for an equation with an
additional quadratic term; they also show that without dispersion all
meromorphic travelling waves are stationary. The elliptic waves of the
Kuramoto--Sivashinsky equation are due to \citet{Kudryashov1990}, and
the condition \(\delta^2=16\mu\nu\) to \citet{Eremenko2006}; it is an
instance of the residue criterion for elliptic solutions
\citep{Hone2005,ConteNg2010}, which Theorem~\ref{thm:C} applies to
\(u^{n-1}\). For \(p=2\), waves outside \(\Wclass\) were
found by \citet{AblowitzZeppetella1979}; see also
\citep[Lemmas~4.4 and~4.5]{ConteNgWu2015}. The estimates of
Section~\ref{sec:analytic} rest on the weighted maximum of Zalcman's
rescaling lemma \citep{Zalcman1975,Zalcman1998}, which Zalcman traces
back to Landau's proof of Bloch's theorem; Lemma~\ref{lem:distance} is
a pointwise counterpart of Clunie's lemma \citep{Laine1993}, and
\citet*{PolacikQuittnerSouplet2007} prove estimates of the same form
for superlinear elliptic equations. For the Painlev\'e transcendents,
Steinmetz \citep{Steinmetz2002,Steinmetz2017,Steinmetz2018} uses
estimates away from the poles and combinations of a solution and its
derivative that are regular at the poles; expressions without poles
built from the first Laurent coefficients are used, with Nevanlinna
theory, by \citet{HalburdWang2015} and \citet{MengHalburd2025}.
Functions of the second kind and Hermite's decomposition are classical
\citep{Hermite1885,AppellLacour1897,DeminaKudryashov2011}.

\emph{New here.} Theorem~\ref{thm:A} for operators of even degree that
are not even, for every \(n\) and for an arbitrary meromorphic
solution; for \(n\ge3\), the single pole modulo twisted periods; the
bound of Theorem~\ref{thm:B} beyond odd \(p\) and even operators of
degree four, with the fact that two waves require an even operator and
\(n\le3\); the formula of Theorem~\ref{thm:C} for \(u^{n-1}\), with the
relation \(u_+^{n-1}-u_-^{n-1}=k\,\mathcal R(P)\) between the limits at
the two ends; and the completeness of the list of Section~\ref{sec:sh}.
The proofs use Cauchy's estimate in place of Nevanlinna theory.

% ======================================================================
\section{The argument on one equation}\label{sec:example}\label{sec:strategy}

We first carry out the whole argument on the equation
\(v''''-sv+v^3/3=0\), that is, \(n=3\) and \(P(D)=D^4-s\); it is the
travelling-wave equation of \(U_t+U^2U_x+U_{xxxxx}=0\) at speed \(s\)
with zero integration constant. There are seven steps. Steps~1--5
are carried out for every operator in
Sections~\ref{sec:local}--\ref{sec:proof}, and Step~7 in
Section~\ref{sec:power}; each of these steps names the
statements that replace it. Step~6 exhibits the two waves of this
equation.

\subsection{Step 1: the Laurent series at a pole}\label{sec:ex-laurent}

By \eqref{eq:balance}, every pole of a wave is double and its leading
coefficient \(c\) satisfies \(c^2=-360\). For either root \(c\), the
function \(u=v/c\) satisfies
\begin{equation}\label{eq:ex-ode}
 u''''=s\,u+120\,u^3 .
\end{equation}
Put a pole at \(0\) and write \(u=z^{-2}(1+b_1z+b_2z^2+\cdots)\). For
every \(j\ge1\), the coefficients of \(z^{j-6}\) on both sides give
\[
 \bigl((j-2)(j-3)(j-4)(j-5)-360\bigr)\,b_j
 =\text{a polynomial in \(s\) and \(b_1,\dots,b_{j-1}\)} .
\]
The product \((j-2)(j-3)(j-4)(j-5)\) is \(24\) or \(0\) for \(j\le5\)
and increases for \(j\ge6\); it equals \(360=3\cdot4\cdot5\cdot6\) only
for \(j=8\). So the factor on the left vanishes only there, and
\(b_1,\dots,b_7\) are determined by \(s\). At \(j=8\) the
coefficient \(b_8\) drops out, and the other terms,
\(-b_4(s+360b_4)\), cancel because \(b_4=-s/360\). So \(b_8\) can be
any number \(h\), and the later coefficients are determined by \(s\)
and \(h\). The result is the series
\[
 f(z;h)=z^{-2}-\frac{s}{360}\,z^{2}+h\,z^{6}+O(z^{10}),
\]
and at every pole \(z_i\) of every wave
\begin{equation}\label{eq:ex-germs}
 v(z_i+z)=\pm c\,f(z;h_i)\qquad\text{for some number }h_i .
\end{equation}
We call \(h_i\) the \emph{free coefficient} of the wave at \(z_i\). It
is all that the equation leaves open near a pole, and it enters the
series late, at \(z^6\). The rest of the argument is about its values.

\emph{In general} (Section~\ref{sec:local}): for odd \(p\) there is
no free coefficient; for even \(p\) there is one, at \(z^{nq}\)
(Lemma~\ref{lem:resonance}, Corollary~\ref{cor:local}).

\subsection{Step 2: a wave is large only near its poles}\label{sec:ex-estimate}

We quote one estimate, proved in Section~\ref{sec:analytic}: every wave
satisfies
\(|v(z)|\le C\bigl(1+d(z)^{-2}\bigr)\), where \(d(z)\) is the distance
from \(z\) to the nearest pole and \(C\) depends only on \(s\)
(Lemma~\ref{lem:distance}). The reason is that on a disk without poles
where \(|v|\) is large, Cauchy's estimate makes \(v''''\) too small to
balance \(v^3/3\). Two consequences are used.
\begin{enumerate}[label=\textup{(\alph*)},leftmargin=2.2em]
\item If the poles of \(v\) form finitely many cosets \(z_i+\Gamma\) of
its group of twisted periods, then \(v\) is elliptic, or a rational
function of an exponential that is bounded at both ends, or a rational
function bounded at infinity, according as \(\Gamma\) is a lattice, a
group \(\tau\ZZ\) with \(\tau\ne0\) (rank one), or \(\{0\}\)
(Lemma~\ref{lem:germs}(c)).
\item A polynomial in \(v,v',v'',\dots\) that has no poles and is
bounded on the set of poles of \(v\) is constant
(Lemma~\ref{lem:finite-values}).
\end{enumerate}

\subsection{Step 3: equal series give a twisted period}\label{sec:ex-global}

If two poles \(z_1,z_2\) of a wave have the same free coefficient, then
\(v(z_2+z)=\pm v(z_1+z)\) near \(z=0\) by \eqref{eq:ex-germs}, hence
for all \(z\) by the identity theorem: \(z_2-z_1\) is a twisted period.
For ordinary periods this argument is due to
\citet{Eremenko1982,Eremenko2006}. So the poles of a wave that carry
one value of \(h\) form a single coset \(z_1+\Gamma\). In the same way, two waves that have poles with
the same value of \(h\) differ by a translation and a sign.

Hence a wave whose poles all carry the same value of \(h\) has one pole
modulo \(\Gamma\), and by Step~2(a) it is elliptic, rational in an
exponential, or rational. It remains to find the possible values of
\(h\) (Step~4) and to show that one wave carries only one of them
(Step~5).

\emph{In general} (Section~\ref{sec:global}):
Corollary~\ref{cor:one-germ} for equal series, and
Proposition~\ref{prop:spaces} for the forms (i)--(iii) of
Theorem~\ref{thm:A}. For odd \(p\) there is no free coefficient, and
this step already proves Theorems~\ref{thm:A} and~\ref{thm:B}
(Section~\ref{sec:odd}).

\subsection{Step 4: an expression without poles}\label{sec:selector-example}

The values of \(h\) are found from a polynomial in \(u\) and its
derivatives that has no poles. The combination
\(uu''-\tfrac32(u')^2\) vanishes for \(u=z^{-2}\). On the series \(f\)
it therefore starts with the next coefficient:
\(ff''-\tfrac32(f')^2=-\tfrac{s}{18}\,z^{-2}+O(z^2)\). We want to
cancel this pole by a multiple of a derivative of \(u\), at every pole
of \(u\) and for either sign in \eqref{eq:ex-germs}. A derivative of
\(u\) changes sign with \(u\), and the quadratic combination does not.
So we first multiply by \(f\), which makes the expression odd in \(u\),
and then compare with \(f''=6z^{-4}+O(1)\):
\[
 f^2f''-\tfrac32\,f\,(f')^2=-\frac{s}{18}\,z^{-4}+O(1),\qquad
 \frac{s}{108}\,f''=\frac{s}{18}\,z^{-4}+O(1).
\]
Neither principal part contains \(h\), and they cancel in the sum. So
\[
 \Psi=u^2u''-\tfrac32\,u\,(u')^2+\frac{s}{108}\,u''
\]
has no pole at a pole of \(u\), whichever sign occurs in
\eqref{eq:ex-germs}: it is an entire function. With more terms of \(f\)
one finds, at a pole \(z_0\) where \(u(z_0+z)=\pm f(z;h)\),
\[
 \Psi(z_0+z)=\pm\Bigl[\,72\,h+\frac{7s^2}{97200}
 +\frac{7\,(388800\,h-s^2)(6998400\,h+7s^2)}{365316480000}\,z^{8}
 +O(z^{12})\Bigr].
\]
The value of \(\Psi\) at a pole contains the unknown \(h\), but
\(\Psi'\) vanishes at every pole of \(u\). By Step~2(b), \(\Psi'\) is
constant, hence zero, and \(\Psi\) is constant. The coefficient of
\(z^8\) must therefore vanish:
\[
 h=\frac{s^2}{388800}\qquad\text{or}\qquad h=-\frac{7s^2}{6998400}.
\]
Every pole of every wave carries one of these two values of the free
coefficient.

\emph{In general} (Section~\ref{sec:selectors}): two expressions in
\(u\) whose principal parts do not contain \(h\) are combined so that
their poles cancel. This gives an equation for \(h\) of degree one if
\(n\ge4\) (Proposition~\ref{prop:affine}) and of degree two if
\(n=2,3\) (Proposition~\ref{prop:quadratic}).

\subsection{Step 5: one wave cannot carry both values}\label{sec:ex-two-values}

Suppose a wave \(v\) has poles \(z_1\) and \(z_2=z_1+a\) with free
coefficients \(h_1\ne h_2\). By Step~4 every pole of \(v\) carries
\(h_1\) or \(h_2\), so by Step~3 its poles form the two cosets
\(z_1+\Gamma\) and \(z_2+\Gamma\), and Step~2(a) applies. Choose the
sign in \(W(x)=v(z_1+a+x)\mp v(z_1+x)\) so that the leading terms at
\(x=0\) cancel. The two series \eqref{eq:ex-germs} agree up to the term
\(h\,z^6\), so \(W(x)=\pm c\,(h_2-h_1)\,x^6+\cdots\). Thus \(W\) has a
zero of order \(6\) at \(0\), and at every point of \(\Gamma\), because
\(W\) has the same multipliers as \(v\). Its poles lie in the cosets
\(\pm a+\Gamma\) and are at most double. If \(\Gamma\) is a lattice,
then \(W^2\) is a function on the torus \(\CC/\Gamma\) with a zero of
order \(12\) and poles of total order at most \(4+4\). This is
impossible: a nonzero meromorphic function on a compact Riemann surface
has as many zeros as poles. If \(\Gamma\) has rank one or is \(\{0\}\),
then \(v\) is bounded at the two ends of the cylinder \(\CC/\Gamma\),
or at infinity, by Step~2(a); so \(W^2\) extends to a sphere, and the
same count applies. For an even operator such as \(D^4-s\), the
conclusion also follows from the first integral obtained with the
factor \(v'\): it takes one value on the wave, and this value is affine
in \(h\) with a nonzero coefficient (Lemma~\ref{lem:energy} and
Theorem~\ref{thm:general}; for \(D^p+a_0\) see
\citet*[Lemma~4]{EremenkoLiaoNg2009}). The count of zeros and poles
works for every operator.

So all poles of a wave carry the same value of \(h\). By Step~3 the
wave has one pole modulo \(\Gamma\) and is elliptic, rational in an
exponential, or rational: this is the first assertion of
Theorem~\ref{thm:A} for \eqref{eq:ex-ode}, and the forms (i)--(iii)
follow as in Proposition~\ref{prop:spaces}. Each of the two values of \(h\) belongs to at most
one wave up to translation and sign: this is the bound of
Theorem~\ref{thm:B}.

\emph{In general} (Section~\ref{sec:proof}): Lemma~\ref{lem:two-germs}
for two values in one wave, where \(n=2\) needs the odd part of \(W\).
If \(n=2,3\) and \(P\) is not even, a further equation of degree one
leaves a single value of \(h\) (Proposition~\ref{prop:not-even}).

\subsection{Step 6: the two waves}\label{sec:two-waves}

Both values of \(h\) occur. For \(s=108\) equation \eqref{eq:ex-ode}
is
\begin{equation}\label{eq:two-wave-ode}
 u''''=108\,u+120\,u^3 ,
\end{equation}
and two solutions are
\begin{equation}\label{eq:two-waves}
 u_1=\wp(z;-6,0),\qquad
 u_2=-y',\quad y^2=\wp(z;4,0),\quad y=z^{-1}+O(z^3).
\end{equation}
Here \(\wp(z;g_2,g_3)\) is the Weierstrass function with
\((\wp')^2=4\wp^3-g_2\wp-g_3\); both discriminants \(g_2^3-27g_3^2\)
are nonzero, so these are elliptic functions and not degenerate ones.
The Weierstrass identities give \(u_1''=6u_1^2+3\) and then
\(u_1''''=120u_1^3+108u_1\). The function \(\wp(z;4,0)\) has double
poles at the points of its period lattice \(\Gamma_0\) and double
zeros at the half-periods where \(\wp=0\), so it has a meromorphic
square root \(y\) on \(\CC\). From \((y')^2=y^4-1\) and \(y''=2y^3\)
we get \(u_2^2=y^4-1\) and
\(u_2''''=(120y^4-12)\,u_2=120u_2^3+108u_2\).

The solution \(u_1\) has trivial character; \(u_2\) is twisted. Indeed,
\(y\) has a single simple pole in a period cell of \(\Gamma_0\), so it
is not periodic for \(\Gamma_0\): its multipliers under \(\Gamma_0\)
form a character \(\chi\) of order two, and \(y\) is an ordinary
elliptic function for the sublattice \(\ker\chi\) of index two, with
two simple poles in a period cell. The function \(u_2=-y'\) has the same
character: its twisted periods are \(\Gamma=\Gamma_0\) and its
ordinary periods are \(\ker\chi\), as in Theorem~\ref{thm:A}(iii).

The Laurent expansions are
\begin{equation}\label{eq:two-wave-jets}
\begin{aligned}
 u_1&=z^{-2}-\tfrac{3}{10}z^2+\tfrac{3}{100}z^6+O(z^{10}),\\
 u_2&=z^{-2}-\tfrac{3}{10}z^2-\tfrac{7}{600}z^6+O(z^{10}),
\end{aligned}
\end{equation}
so the free coefficients are \(h_1=3/100\) and \(h_2=-7/600\), the two
values of Step~4 for \(s=108\). So \(v_1=c\,u_1\) and \(v_2=c\,u_2\),
with \(c^2=-360\), are two waves that are not related by translation or
phase change. They are all the waves of this equation: a wave has a
pole, by Step~4 its free coefficient there is \(h_1\) or \(h_2\), and
by Step~3 the wave is then \(v_1\) or \(v_2\) up to a translation and a
sign.

If \(u\) solves \eqref{eq:ex-ode}, then \(\lambda^2u(\lambda z)\)
solves it with \(\lambda^4s\) in place of \(s\). So
\(c\,\lambda^2u_{1,2}(\lambda z)\) are the two waves for
\(P(D)=D^4-108\lambda^4\). Every \(s\ne0\) has the form
\(108\lambda^4\); hence \(v''''-sv+v^3/3=0\) has exactly two waves for
every \(s\ne0\), as Theorem~\ref{thm:B} asserts for \(n=3\)
(Section~\ref{sec:n2} treats \(n=2\)). For
\(s=0\) the two values of Step~4 coincide, \(h=0\); the series
\(f(z;0)\) is then \(z^{-2}\), and the only wave is \(c\,z^{-2}\). The
two waves are the solutions (36) and (51) of
\citet{DeminaKudryashovKawahara2010} with \(\beta=0\) and \(C_1=0\),
and the solutions (17) and (22) of \citet*{ConteNgWong2012} with
\(b=s_1=s_3=0\), in the notation of those papers.

\subsection{Step 7: the square of a wave}\label{sec:ex-power}

The series of Step~1 also determines the square of a wave, once
\(\Gamma\) is known. It gives \(u^2=z^{-4}-\tfrac{s}{180}+O(z^4)\) at a
pole of any wave: in the notation of Theorem~\ref{thm:C}, \(\ell_0=1\),
\(\ell_1=\ell_2=\ell_3=0\) and \(\ell_4=-s/180\). The function \(u^2\)
is invariant under \(\Gamma\), because the multipliers of \(u\) are
\(\pm1\), and it has one pole modulo \(\Gamma\). Let \(\Gamma\) be a
lattice and \(\wp=\wp(z;g_2,g_3)\) its Weierstrass function, so that
\(\wp^2=z^{-4}+g_2/10+O(z^2)\). Then \(u^2-\wp^2\) is an elliptic
function without poles, hence a constant:
\[
 u^2=\wp^2-\frac{g_2}{10}-\frac{s}{180}.
\]
So the square of a wave is known as soon as its lattice is known; the
character does not enter. For \(s=108\) the formula gives
\(u_1^2=\wp(z;-6,0)^2\) and \(u_2^2=\wp(z;4,0)^2-1\), in agreement with
Step~6.

\emph{In general} (Section~\ref{sec:power}): Theorem~\ref{thm:power}
gives \(u^{n-1}\) for each kind of \(\Gamma\), and
Corollary~\ref{cor:residue} the conditions on its residue, which is
\(\ell_3=0\) here.

% ======================================================================
\section{The Laurent series at a pole}\label{sec:local}\label{sec:origin}\label{sec:laurent}

This section does Step~1 of Section~\ref{sec:example} for every operator. From now on \(p/(n-1)\)
is a positive integer, as it must be if a wave exists
(Proposition~\ref{prop:existence}). We fix the notation
\[
 m=n-1,\qquad p=mq,\qquad
 K=(-1)^p(q)_p=(-1)^p\,\frac{(nq-1)!}{(q-1)!},
\]
and a number \(c\) with \(c^m=-nK\), as in \eqref{eq:balance}. If a
wave \(v\) has a pole at \(z_0\), a phase change makes its leading
coefficient there equal to \(c\), and then \(u(z)=v(z_0+z)/c\)
satisfies
\begin{equation}\label{eq:normalized}
 P(D)u=K\,u^n,\qquad u=z^{-q}\bigl(1+O(z)\bigr).
\end{equation}
We call such a \(u\) a \emph{normalized wave}; it has its pole at
\(0\).

We write
\[
 P(\rho)=\rho^p+a_{p-1}\rho^{p-1}+\dots+a_1\rho+a_0 .
\]
If \(u\) is a normalized wave for \(P\) and \(\lambda\ne0\), then
\(\lambda^qu(\lambda z)\) is a normalized wave for
\(P_\lambda(\rho)=\lambda^pP(\rho/\lambda)\). This \emph{scaling}
multiplies \(a_{p-k}\) by \(\lambda^k\); accordingly we say that
\(a_{p-k}\) has \emph{weight} \(k\).

Let \(u\) be a normalized wave and write
\(u=z^{-q}\sum_{j\ge0}b_jz^j\) with \(b_0=1\). The coefficient of
\(z^{j-q-p}\) in \(P(D)u-Ku^n\) gives the recurrence
\begin{equation}\label{eq:recurrence}
 I(j)\,b_j=F_j(b_1,\dots,b_{j-1}),\qquad
 I(j)=(j-q)(j-q-1)\cdots(j-q-p+1)-nK .
\end{equation}
The left-hand side collects the terms that are linear in \(b_j\),
coming from \(D^p\) and from \(u^n\). The right-hand side \(F_j\)
collects the rest; it is a polynomial in \(b_1,\dots,b_{j-1}\) and in
the coefficients of \(P\). The same recurrence holds at every pole of
every wave, because \eqref{eq:normalized} does not depend on which root
\(c\) was chosen. Which coefficients it determines depends on the
integer roots of the \emph{indicial polynomial} \(I\).

\begin{lemma}[The resonance]\label{lem:resonance}
If \(p\) is odd, then \(I(j)\ne0\) for every integer \(j\ge1\). If
\(p\) is even, the only positive integer root of \(I\) is
\[
 J=p+2q=(n+1)q .
\]
\end{lemma}

\begin{proof}
Write \(\Pi(j)=(j-q)(j-q-1)\cdots(j-q-p+1)\) for the product in
\(I(j)\). Since \(q+p=nq\), we have \(n(q)_p=(q+1)_p\). For
\(1\le j<q\), all factors of \(\Pi(j)\) are negative, and
\(|\Pi(j)|=(q-j)_p<(q)_p<n(q)_p=|nK|\), so \(I(j)\ne0\). For
\(q\le j<q+p\), one factor
vanishes and \(I(j)=-nK\ne0\). For \(j\ge q+p\), we have
\(\Pi(j)=(j-q-p+1)_p>0\), and this is strictly increasing in \(j\). If
\(p\) is odd, then \(-nK=n(q)_p>0\), so \(I(j)>0\). If \(p\) is even,
then \(nK=(q+1)_p\), and \(I(j)=0\) exactly when \(j-q-p+1=q+1\), that
is, \(j=p+2q\).
\end{proof}

In the language of the Painlev\'e test, the roots of \(I\) are the
resonances; for even \(p\) the only positive integer resonance is \(J\)
\citep{Wu2014}. So for odd \(p\) the
recurrence determines all \(b_j\), and the Laurent series of a
normalized wave is determined by \(P\). For even \(p\) the recurrence determines
\(b_1,\dots,b_{J-1}\); at the index \(J\) it reads
\(F_J(b_1,\dots,b_{J-1})=0\) and does not contain \(b_J\); and it
determines \(b_j\) for \(j>J\) once \(b_J\) is given. For every
operator \(P\) of even degree and every \(h\in\CC\) we define the
series
\begin{equation}\label{eq:family}
 f(z;h)=z^{-q}\sum_{j\ge0}b_j(h)\,z^j,\qquad b_0=1,\quad b_J=h,
\end{equation}
by the recurrence at all indices \(j\ne J\). Then
\(P(D)f-Kf^n=-F_J\,z^{q}\), where the number \(F_J\) depends on \(P\)
and not on \(h\). If \(P\) admits a wave, the wave has a pole
(Proposition~\ref{prop:existence}), and its Laurent series there
satisfies the recurrence at the index \(J\) as well. So \(F_J=0\), and
\(f(z;h)\) is a formal solution of \eqref{eq:normalized} for every
\(h\). (We define the series for every operator, because
Sections~\ref{sec:selectors} and~\ref{sec:proof} treat the quantities
built from it as polynomials in the coefficients of \(P\).) At each
pole \(z_i\) of a wave \(v\) we then have
\begin{equation}\label{eq:pole-germs}
 v(z_i+z)=c\,\omega_i\,f(z;h_i)\qquad
 \text{for some }\omega_i\in\mu_m,\ h_i\in\CC.
\end{equation}
We call \(h_i\) the \emph{free coefficient} of the wave at the pole
\(z_i\), \(\omega_i\) its \emph{phase} there, and \(f(\cdot\,;h_i)\) its
\emph{normalized series} there. For
odd \(p\) we use the same notation without \(h\).

\begin{lemma}[Weights]\label{lem:degree}
Let \(p\) be even. Each \(b_j(h)\) is a polynomial in \(h\) and in the
coefficients of \(P\). It is weighted-homogeneous of weight \(j\) if
each \(a_{p-i}\) is given its weight \(i\) and \(h\) the weight \(J\). In
particular, \(b_j(h)\) has degree at most \(\lfloor j/J\rfloor\) in
\(h\); the part of \(f(z;h)\) of degree \(k\ge1\) in \(h\) is
\(O(z^{kJ-q})\); and the part of degree one is
\(h\,z^{J-q}\bigl(1+O(z)\bigr)\).
\end{lemma}

\begin{proof}
The recurrence gives \(b_j\) as a polynomial, because \(I(j)\) is a
nonzero number for \(j\ne J\). For \(\lambda\ne0\), the series
\(\lambda^qf(\lambda z;h)=z^{-q}\sum_j\lambda^jb_j(h)z^j\) starts with
\(z^{-q}\), has the coefficient \(\lambda^Jh\) at the index \(J\), and
satisfies
\[
 P_\lambda(D)\bigl[\lambda^qf(\lambda z;h)\bigr]
 -K\bigl[\lambda^qf(\lambda z;h)\bigr]^n
 =\lambda^{nq}\bigl(P(D)f-Kf^n\bigr)(\lambda z;h)
 =-\lambda^{J}F_J\,z^{q}.
\]
So it satisfies the recurrence of \(P_\lambda\) at all indices
\(j\ne J\), and it is the series \eqref{eq:family} of the operator
\(P_\lambda\) with free coefficient \(\lambda^Jh\). The coefficients of
\(P_\lambda\) are \(\lambda^ia_{p-i}\), which proves the homogeneity.
All weights are positive, so a monomial of \(b_j\) that contains
\(h^k\) has \(kJ\le j\). The last two assertions follow, because
\(b_j\) does not depend on \(h\) for \(j<J\) and \(b_J=h\).
\end{proof}

\begin{corollary}\label{cor:local}
\textup{(a)} The principal part and the constant term of a normalized
series are determined by \(P\).
\textup{(b)} If \(z_1,z_2\) are poles of a wave \(v\), written as in
\eqref{eq:pole-germs}, then
\[
 v(z_2+z)-\frac{\omega_2}{\omega_1}\,v(z_1+z)=O(z^{nq}).
\]
\end{corollary}

\begin{proof}
(a) These are the coefficients \(b_j\) with \(j\le q\), and \(q<J\)
when \(p\) is even. (b) The two
normalized series have the same coefficients \(b_0,\dots,b_{J-1}\), so
the difference is \(c\,\omega_2\bigl(f(z;h_2)-f(z;h_1)\bigr)
=O(z^{J-q})\), and \(J-q=nq\).
\end{proof}

Thus the free coefficient appears only at the order \(z^{nq}\), long
after the pole; Sections~\ref{sec:selectors} and~\ref{sec:proof} rest
on this.

% ======================================================================
\section{A wave is large only near its poles}\label{sec:analytic}

This section proves the estimate of Step~2 of
Section~\ref{sec:example} and its consequence~(b):
Lemma~\ref{lem:distance} bounds a solution in terms of the distance to
its poles, and Lemma~\ref{lem:finite-values} says that an expression
without poles is constant if it is bounded on the poles. The first
is used in Section~\ref{sec:global} to place a wave in \(\Wclass\); the
second, in Section~\ref{sec:selectors}, to obtain the equations for the
free coefficient. They are the only analytic estimates in
the paper, and they do not use the form of the nonlinearity. We prove them for the solutions, meromorphic in \(\CC\),
of
\begin{equation}\label{eq:general-eq}
 P(D)v+B(v)=0,
\end{equation}
where \(P(D)=\sum_{j=0}^pa_jD^j\) is any operator with constant
coefficients, \(a_p\ne0\) and \(p\ge1\), and \(B\) is any polynomial of
degree \(n\ge2\); here \(q=p/(n-1)\) need not be an integer. This
generality is used in Proposition~\ref{prop:existence} and in
Appendix~\ref{app:general}; for Theorems~\ref{thm:A}--\ref{thm:C} the
reader may take \(B(v)=v^n/n\). For a solution \(v\), let \(S\) be its
set of poles and \(d(z)=\dist(z,S)\); if \(S\) is empty, we put
\(d\equiv\infty\) and \(d^{-q}=0\).

Both lemmas rest on the maximum of a weighted modulus, the device of
Zalcman's rescaling lemma \citep{Zalcman1975,Zalcman1998}. In a given
disk it produces a point \(z_*\) at which the modulus of a holomorphic
function \(w\), weighted by a power of the distance to the boundary, is
at least as large as at the centre, and such that \(w\) stays close to
\(w(z_*)\) on a disk about \(z_*\) whose radius is proportional to the
distance from \(z_*\) to the boundary.

\begin{lemma}[A weighted maximum]\label{lem:weighted}
Let \(w\) be holomorphic on a neighbourhood of the closed disk
\(\overline D(z_0,R)\) and not identically zero on it, and let
\(\theta>0\). There is a point \(z_*\) in the open disk such that, with
\(\delta=R-|z_*-z_0|\),
\[
 R^\theta|w(z_0)|\le \delta^\theta|w(z_*)|
 \qquad\text{and}\qquad
 \max_{|z-z_*|\le \delta/2}|w(z)|\le 2^\theta|w(z_*)|.
\]
Moreover, with \(c_\theta=2^{-\theta-3}\),
\[
 |w(z)-w(z_*)|\le\tfrac13\,|w(z_*)|
 \qquad\text{for } |z-z_*|\le c_\theta\delta .
\]
In particular, \(w\) has no zero in \(\overline D(z_*,c_\theta\delta)\).
\end{lemma}

\begin{proof}
Let \(z_*\) be a point where \((R-|z-z_0|)^\theta|w(z)|\) attains its
maximum on the closed disk. The maximum is positive and the weight
vanishes on the boundary, so \(z_*\) is an interior point, and the
first inequality is the comparison with the value at \(z_0\). On
\(\overline D(z_*,\delta/2)\) the weight is at least
\((\delta/2)^\theta\), which gives the second inequality. Put
\(M=|w(z_*)|>0\). Cauchy's estimate on \(\overline D(z_*,\delta/2)\)
bounds the Taylor coefficients of \(w\) at \(z_*\), and gives, for
\(|t|<\delta/2\),
\[
 |w(z_*+t)-w(z_*)|\le 2^\theta M\sum_{j\ge1}\Bigl(\frac{2|t|}{\delta}\Bigr)^j .
\]
For \(|t|\le c_\theta\delta\) the right-hand side is at most \(M/3\).
\end{proof}

\begin{lemma}[Distance to the poles]\label{lem:distance}
There are constants \(C_j\), depending only on \(P\), \(B\) and \(j\),
such that every meromorphic solution of \eqref{eq:general-eq} satisfies,
for \(z\notin S\),
\begin{equation}\label{eq:distance}
 |v(z)|\le C_0\bigl(1+d(z)^{-q}\bigr),\qquad
 |v^{(j)}(z)|\le C_j\bigl(1+d(z)^{-q-j}\bigr)\quad(j\ge1).
\end{equation}
If \(S\) is finite, then \(v\) is a rational function bounded at
infinity.
\end{lemma}

\begin{proof}
For \(v\equiv0\) there is nothing to prove. Choose \(C_B>0\) and
\(V_0\ge1\) such that \(|B(x)|\ge C_B|x|^n\) for \(|x|\ge V_0\).
Fix \(z_0\notin S\) and \(0<R<d(z_0)\). Apply
Lemma~\ref{lem:weighted} to \(v\) on \(\overline D(z_0,R)\) with
\(\theta=q\), and put \(V=|v(z_*)|\) and \(r=\delta/2\). Suppose
\(V\ge V_0\). Since \(|v|\le2^qV\) on \(\overline D(z_*,r)\), Cauchy's
estimate gives \(|v^{(j)}(z_*)|\le j!\,2^qVr^{-j}\), and the equation
gives
\[
 C_BV^{n}\le|B(v(z_*))|=|P(D)v(z_*)|
 \le 2^qV\sum_{j=0}^{p}|a_j|\,j!\,r^{-j}.
\]
Hence \(V^{n-1}\le C'\max(1,r^{-p})\), and so
\(V\le C(1+r^{-q})\), with \(C\) depending only on \(P\) and \(B\).
If \(V<V_0\), the same bound holds as soon as \(C\ge V_0\). The first inequality of
Lemma~\ref{lem:weighted} now gives
\[
 R^q|v(z_0)|\le \delta^qV\le C(\delta^q+2^q)\le C(R^q+2^q).
\]
Letting \(R\to d(z_0)\), or \(R\to\infty\) if \(S\) is empty, proves the
first estimate.

For the derivatives at \(z\), use the disk of radius
\(r_1=\tfrac12\min(1,d(z))\). On this disk the distance to \(S\) is at
least \(r_1\), so the first estimate bounds \(v\) there by
\(C_0(1+r_1^{-q})\), and Cauchy's estimate bounds \(v^{(j)}(z)\) by
\(j!\,r_1^{-j}C_0(1+r_1^{-q})\). This is the second estimate.

If \(S\) is finite, then \(d(z)\to\infty\) as \(z\to\infty\), so \(v\)
is bounded near infinity. It extends holomorphically to \(\infty\), and
a meromorphic function on the Riemann sphere is rational.
\end{proof}

\begin{proposition}[Solutions without poles]\label{prop:existence}
Every entire solution of \eqref{eq:general-eq} is constant. If
\(p/(n-1)\) is not a positive integer, every meromorphic solution of
\eqref{eq:general-eq} is constant.
\end{proposition}

\begin{proof}
An entire solution is bounded by Lemma~\ref{lem:distance}, hence
constant by Liouville's theorem. At a pole of any solution the balance
of \(v^{(p)}\) with the term of degree \(n\) in \(B(v)\) gives the
order \(p/(n-1)\). If this is not a positive integer there is no pole,
and the solution is entire.
\end{proof}

For example, the travelling waves of the generalized Bretherton
equation \(U_{tt}+\alpha U_{xx}+\beta U_{xxxx}+\delta U+\gamma U^n=0\)
with \(\beta\gamma\ne0\) satisfy, after a rescaling, \eqref{eq:main}
with an even operator of degree four, so waves exist only for
\(n\in\{2,3,5\}\) (for elliptic waves and nonzero coefficients of
\(D^2\) and \(D^0\) in \(P\), see
\citep{KudryashovSinelshchikovDemina2011}).

\begin{lemma}[Bounded on the poles]\label{lem:finite-values}
Let \(v\) be a nonconstant meromorphic solution of
\eqref{eq:general-eq}, and let \(\Phi\) be a polynomial in
\(v,v',v'',\dots\) with constant coefficients. If \(\Phi\) has no poles
and is bounded on \(S\), then \(\Phi\) is constant.
\end{lemma}

\begin{proof}
We may assume \(\Phi\not\equiv0\). Put \(M=\sup_S|\Phi|\). By
\eqref{eq:distance} there are constants \(C,N>0\) such that
\(|\Phi(z)|\le C\bigl(1+d(z)^{-N}\bigr)\) for \(z\notin S\). Given any
disk \(\overline D(z_0,R)\), apply Lemma~\ref{lem:weighted} to the
entire function \(\Phi\) with \(\theta=N\). If the disk
\(\overline D(z_*,c_N\delta)\) contains a pole \(z_i\) of \(v\), then
\(|\Phi(z_*)|\le\tfrac32|\Phi(z_i)|\le\tfrac32M\). If it contains no
pole, then \(d(z_*)\ge c_N\delta\) and
\(|\Phi(z_*)|\le C\bigl(1+(c_N\delta)^{-N}\bigr)\). In both cases
\[
 R^N|\Phi(z_0)|\le \delta^N|\Phi(z_*)|
 \le\max\bigl(\tfrac32M,\,C\bigr)\,R^N+C\,c_N^{-N}.
\]
Letting \(R\to\infty\) gives \(|\Phi(z_0)|\le\max(\tfrac32M,C)\) for
every \(z_0\). So \(\Phi\) is bounded, and constant by Liouville's
theorem.
\end{proof}

In the applications \(\Phi\) takes only finitely many values on \(S\).
The hypothesis on the values at the poles cannot be dropped; an example
with \(p=2\) is given in Appendix~\ref{sec:low-order}.

% ======================================================================
\section{From one Laurent series to the global structure}\label{sec:global}

This section does Step~3 of Section~\ref{sec:example} in general.

\subsection{Equal series give twisted periods}

From here on we return to waves. Recall that \(\Gamma\) is the group of
twisted periods of a wave, \(\chi\colon\Gamma\to\mu_m\) its character
(Definition~\ref{def:twisted}), and \(S\) its set of poles. A discrete
subgroup of \(\CC\) is \(\{0\}\), a group \(\tau\ZZ\) with
\(\tau\ne0\), or a lattice; we say that it has \emph{rank} zero, one or
two.

\begin{lemma}[Equal series and twisted periods]\label{lem:germs}
Let \(v\) be a wave.
\begin{enumerate}[label=\textup{(\alph*)},leftmargin=2.2em]
\item Two poles \(z_1,z_2\) of \(v\) satisfy \(z_2-z_1\in\Gamma\) if
and only if \(v(z_2+z)=\omega\,v(z_1+z)\) near \(z=0\) for some
\(\omega\in\mu_m\).
\item \(\Gamma\) is a discrete subgroup of \(\CC\), and the group
\(\ker\chi\) of ordinary periods has finite index in it.
\item Suppose the poles of \(v\) lie in finitely many cosets of
\(\Gamma\). If \(\Gamma=\{0\}\), then \(v\) is rational and bounded
at infinity. If \(\Gamma\) has rank one and \(\tau\) generates
\(\ker\chi\), then \(v\) is a rational function of
\(e^{2\pi iz/\tau}\) that is finite at \(0\) and at \(\infty\). If
\(\Gamma\) is a lattice, then \(v\) is elliptic.
\end{enumerate}
\end{lemma}

\begin{proof}
(a) If \(\gamma=z_2-z_1\in\Gamma\), then
\(v(z_2+z)=v(z_1+z+\gamma)=\chi(\gamma)v(z_1+z)\). Conversely, if the
two Laurent series agree up to the factor \(\omega\), then
\(v(z+z_2-z_1)=\omega v(z)\) near \(z_1\), hence everywhere by the
identity theorem.

(b) Suppose distinct elements \(\gamma_i\in\Gamma\) tend to \(0\).
Passing to a subsequence, they have a common multiplier \(\omega\).
Letting \(i\to\infty\) in \(v(z+\gamma_i)=\omega v(z)\) at a point
where \(v\) is holomorphic and nonzero gives \(\omega=1\). Then
\((v(z+\gamma_i)-v(z))/\gamma_i=0\) gives \(v'\equiv0\), which is
impossible. So \(\Gamma\) is discrete. The kernel of \(\chi\) has
finite index because \(\mu_m\) is finite.

(c) If \(\Gamma=\{0\}\), then \(v\) has finitely many poles, and
Lemma~\ref{lem:distance} applies. If \(\Gamma\) is a lattice, so is
\(\ker\chi\), and \(v\) is elliptic. Suppose \(\Gamma\) has rank one.
Each coset of \(\Gamma\) is a finite union of cosets of
\(\ker\chi=\tau\ZZ\), so \(S=\bigcup_{\nu=1}^{N}(a_\nu+\tau\ZZ)\). With
\(y=\operatorname{Im}(z/\tau)\) we have
\[
 d(z)\ge|\tau|\min_\nu\bigl|y-\operatorname{Im}(a_\nu/\tau)\bigr|
 \longrightarrow\infty\qquad(y\to\pm\infty),
\]
uniformly in \(\operatorname{Re}(z/\tau)\). By
Lemma~\ref{lem:distance}, \(v\) is bounded for large \(|y|\). Since
\(\tau\) is a period, \(v\) is a function of \(t=e^{2\pi iz/\tau}\),
meromorphic on \(\CC^*\) with finitely many poles. As
\(|t|=e^{-2\pi y}\), it is bounded near \(t=0\) and \(t=\infty\). These
singularities are therefore removable, and \(v\) is rational in \(t\).
\end{proof}

All leading coefficients of a wave differ by factors in \(\mu_m\). So
part (a) says: two poles lie in the same coset of \(\Gamma\) exactly
when their normalized series coincide.

\begin{corollary}[One series, one orbit]\label{cor:one-germ}
\textup{(a)} If a wave has the same normalized series at all its poles
(for even \(p\), the same free coefficient at all its poles), then it
has exactly one pole modulo \(\Gamma\) and belongs to \(\Wclass\).

\textup{(b)} If a pole of one wave and a pole of another wave carry the
same normalized series, the two waves coincide after a translation and a
phase change.
\end{corollary}

\begin{proof}
(a) By Lemma~\ref{lem:germs}(a) all poles lie in one coset of
\(\Gamma\), and part (c) applies. For even \(p\) the free coefficient
determines the normalized series. (b) If
\(v_1(z_1+z)=c\,\omega_1f(z)\) and \(v_2(z_2+z)=c\,\omega_2f(z)\) near
\(z=0\) with the same series \(f\), then
\(v_2(z)=(\omega_2/\omega_1)\,v_1(z+z_1-z_2)\) near \(z_2\), hence
everywhere by the identity theorem.
\end{proof}

\subsection{The one-pole spaces}\label{sec:spaces}

We now describe the functions with one pole modulo \(\Gamma\),
according to the rank of \(\Gamma\). For \(\Gamma=(2\pi i/k)\ZZ\) we
use the functions \(g_r\) of Theorem~\ref{thm:A},
\begin{equation}\label{eq:generator}
 g_r(z)=\frac{k\,e^{rkz/m}}{e^{kz}-1},\qquad 1\le r\le m,\quad k\ne0.
\end{equation}
Recall that \(g_r\) has simple poles exactly at the points of
\(\Gamma\), with residue \(1\) at \(0\), and the multiplier
\(e^{2\pi ir/m}\) under \(z\mapsto z+2\pi i/k\). When \(r<m\), \(g_r\) tends to \(0\) at both
ends, while \(g_m=G\) tends to \(k\) at the end
\(\operatorname{Re}(kz)\to+\infty\) and to \(0\) at the other end.
Replacing \(k\) by \(-k\), which gives the same \(\Gamma\), turns
\(g_r\) into \(g_{m-r}\) for \(r<m\), and \(g_m\) into \(g_m-k\); so
the form \eqref{eq:periodic-form} does not depend on the choice between
\(k\) and \(-k\).

For a lattice \(\Gamma\), a character \(\chi\colon\Gamma\to\CC^*\) of
finite order and \(j\ge0\), let \(V_j^\chi\) be the space of
meromorphic functions \(w\) with \(w(z+\gamma)=\chi(\gamma)w(z)\) for
\(\gamma\in\Gamma\), whose poles lie in \(\Gamma\) and have order at
most \(j\). These are elliptic functions of the second kind with
multipliers of finite order. The following description of them is
classical \citep{Hermite1885,AppellLacour1897}. It makes
Theorem~\ref{thm:A}(iii) explicit and is not used in the proofs of the
three theorems; its proof, a construction with the Weierstrass sigma
function, is in Appendix~\ref{app:second-kind}.

\begin{lemma}[Functions of the second kind]\label{lem:second-kind}
Let \(\chi\) be a character of finite order of the lattice \(\Gamma\).
If \(\chi\) is not trivial, then \(V_0^\chi=0\) and
\(\dim V_j^\chi=j\) for \(j\ge1\); there is a unique
\(g_\chi\in V_1^\chi\) with residue \(1\) at \(0\), and
\(g_\chi,Dg_\chi,\dots,D^{j-1}g_\chi\) is a basis of \(V_j^\chi\). If
\(\chi\) is trivial, then \(V_0^\chi=V_1^\chi=\CC\), and for \(j\ge2\)
a basis of \(V_j^\chi\) is \(1,\wp,D\wp,\dots,D^{j-2}\wp\).
\end{lemma}

For \(n=3\) the functions \(g_\chi\) are, up to scale, Jacobi elliptic
functions (Appendix~\ref{app:second-kind}); the function \(y\) of
Section~\ref{sec:two-waves} is one of them.

\begin{proposition}[The one-pole spaces]\label{prop:spaces}
Let \(v\) be a wave with exactly one pole modulo \(\Gamma\), placed at
\(0\). Then \(v\) has the form \textup{(i)}, \textup{(ii)} or
\textup{(iii)} of Theorem~\ref{thm:A}, according to the rank of
\(\Gamma\). In case \textup{(iii)}, \(v\in V_q^\chi\).
\end{proposition}

\begin{proof}
By Lemma~\ref{lem:germs}(c), applied with a single coset, \(v\) is
rational and bounded at infinity, or rational in an exponential and
bounded at both ends, or elliptic.

\emph{Rank zero.} Then \(0\) is the only pole, so
\(v=\beta+\sum_{j=1}^{q}d_jz^{-j}\) with \(d_q\ne0\). Let \(j_0\) be
the smallest index with \(d_{j_0}\ne0\). At infinity, the coefficients
of \(z^0\) and of \(z^{-j_0}\) in \(P(D)v+v^n/n\) are
\(P(0)\beta+\beta^n/n\) and \((P(0)+\beta^m)\,d_{j_0}\). Both vanish, so
\(P(0)=-\beta^m\) and then \((\tfrac1n-1)\beta^n=0\). Hence
\(\beta=0\).

\emph{Rank one.} Let \(\Gamma=(2\pi i/k)\ZZ\) and let
\(e^{2\pi ir/m}\) be the multiplier of \(v\) under \(z\mapsto z+2\pi
i/k\). The functions \(D^jg_r\), \(0\le j<q\), have this multiplier,
their principal parts at \(0\) are \((-1)^jj!\,z^{-j-1}\), and they are
bounded at both ends. Choose \(d_0,\dots,d_{q-1}\) so that the
difference \(v-\sum_jd_jD^jg_r\) is holomorphic at \(0\); then
\(d_{q-1}\ne0\), because the pole of \(v\) has order \(q\). The
difference has the same multiplier as \(v\), so it is holomorphic at
every point of \(\Gamma\), hence entire. It is also bounded at both
ends. As a function of \(e^{2\pi iz/\tau}\), where \(\tau\) is an
ordinary period, it is therefore rational, without poles on \(\CC^*\)
and bounded at \(0\) and \(\infty\), that is, a constant. A nonzero
constant has trivial multiplier, so the difference is zero if \(r<m\),
and it is a constant \(\beta\) if \(r=m\).

\emph{Rank two.} This is the definition of \(V_q^\chi\), with the
character \(\chi\) of \(v\).
\end{proof}

By Lemma~\ref{lem:second-kind}, a twisted elliptic wave is thus a
combination of \(g_\chi,Dg_\chi,\dots,D^{q-1}g_\chi\), in analogy with
\eqref{eq:periodic-form}, and an untwisted one, which requires
\(q\ge2\), is a combination of \(1,\wp,D\wp,\dots,D^{q-2}\wp\).

\subsection{Odd degree}\label{sec:odd}

Let \(p\) be odd. By Lemma~\ref{lem:resonance} all poles of all waves
carry the same normalized series. By Corollary~\ref{cor:one-germ}(a)
every wave has exactly one pole modulo \(\Gamma\), and
Proposition~\ref{prop:spaces} gives its form; by
Corollary~\ref{cor:one-germ}(b) any two waves coincide after a
translation and a phase change. This proves Theorems~\ref{thm:A}
and~\ref{thm:B} for odd \(p\).

For even \(p\) it remains to show\label{rem:remains} that the free
coefficient takes the same value at all poles of a wave, and at most
two values on all waves of one equation, only one when \(n\ge4\) or
\(P\) is not even.

% ======================================================================
\section{Equations for the free coefficient}\label{sec:selectors}

This section does Step~4 of Section~\ref{sec:example} for every operator of even degree: it
constructs expressions like the \(\Psi\) of
Section~\ref{sec:selector-example}, which have no poles and whose
Taylor coefficients at a pole give equations for \(h\). Throughout,
\(p\) is even, so \(K=(q)_p>0\) and \(nK=(q+1)_p\). Recall that the
series \(f(z;h)\) of \eqref{eq:family} is defined for every operator
\(P\) of even degree, with coefficients that are polynomials in \(h\)
and in the coefficients of \(P\). We write \([h^k]\) and \([z^j]\) for
the coefficient of \(h^k\) and of \(z^j\) in such a series.
Table~\ref{tab:notation} in Appendix~\ref{app:notation} lists the
notation.

\subsection{Functions without poles}

In Section~\ref{sec:selector-example} the principal part of
\(u^2u''-\tfrac32u(u')^2\) at a pole was cancelled by a multiple of
\(u''\), that is, by a differential operator applied to \(u\). For a
general operator the principal part of \(u\) at a pole has \(q\) terms,
which depend on \(P\), and one cannot expect the pole of a second
expression to be cancelled by derivatives of \(u\) alone. Instead we
apply a differential operator to each of two expressions, so that
their principal parts cancel against each other.

For a Laurent series \(\alpha\) at \(0\) with principal part
\(\sum_{r=1}^{R}\alpha_{-r}z^{-r}\), define the \emph{principal-part
polynomial}
\[
 \hatt\alpha(\rho)=\sum_{r=1}^{R}\alpha_{-r}\,\frac{(-\rho)^{r-1}}{(r-1)!} .
\]
Since \(D^{r-1}(1/z)=(-1)^{r-1}(r-1)!\,z^{-r}\), the principal part of
\(\alpha\) is \(\hatt\alpha(D)(1/z)\). Differentiation maps principal
parts to principal parts and holomorphic functions to holomorphic
functions, so for every polynomial \(L\) the principal part of
\(L(D)\alpha\) is \(L(D)\hatt\alpha(D)(1/z)\); in particular
\(\hatt{D\alpha}=\rho\,\hatt\alpha\).

\begin{lemma}[Cancelling principal parts]\label{lem:bracket}
For Laurent series \(\alpha\) and \(\beta\) at \(0\) with finite
principal parts, the series
\(\hatt\alpha(D)\,\beta-\hatt\beta(D)\,\alpha\) is holomorphic at
\(0\).
\end{lemma}

\begin{proof}
Its principal part is
\((\hatt\alpha\hatt\beta-\hatt\beta\hatt\alpha)(D)(1/z)=0\).
\end{proof}

For convergent series the lemma has a transparent reason. Let
\(\alpha\) and \(\beta\) be holomorphic in \(0<|\zeta|<2\delta\). For
\(0<|z|<\delta\) the residues at \(\zeta=0\) and \(\zeta=z\) give
\[
 \hatt\alpha(D)\,\beta(z)-\hatt\beta(D)\,\alpha(z)
 =\frac{1}{2\pi i}\oint_{|\zeta|=\delta}\alpha(\zeta)\,\beta(z-\zeta)\,d\zeta ,
\]
and the integral is holomorphic in \(|z|<\delta\).

The function \(\Psi\) of Section~\ref{sec:selector-example} arises from
the lemma. Take \(E=u\) and \(F=u^2u''-\tfrac32u(u')^2\). On the series
\(f(z;h)\) of that example, the principal-part polynomials of \(E\) and
\(F\) are \(\hatt E(\rho)=-\rho\) and
\(\hatt F(\rho)=\tfrac{s}{108}\rho^3\), and
\[
 \hatt E(D)\,F-\hatt F(D)\,E=-D\Bigl(F+\frac{s}{108}\,u''\Bigr)=-\Psi' .
\]

In general we need two expressions in \(u\) with two properties. Their
principal parts on the series \(f(z;h)\) must not depend on \(h\); then
\(\hatt E\) and \(\hatt F\) are computed from \(P\). And the two
expressions must transform in the same way under phase changes; then at
a pole with phase \(\omega_i\) the principal parts of \(E(u)\) and
\(F(u)\) are those of \(E(f)\) and \(F(f)\) times the same power of
\(\omega_i\). So one differential polynomial in \(u\) is holomorphic
at all poles of the wave, and its expansions at the poles differ only
by roots of unity and by the value of \(h\). Thus let \(E\) and \(F\)
be polynomials in \(u,u',u'',\dots\) with constant coefficients such
that:
\begin{enumerate}[label=\textup{(H\arabic*)},leftmargin=2.8em]
\item for some integer \(e\), \(E(\omega u)=\omega^eE(u)\) and
\(F(\omega u)=\omega^eF(u)\) for all \(\omega\in\mu_m\);
\item the principal parts at \(0\) of \(E(f(\cdot\,;h))\) and
\(F(f(\cdot\,;h))\) do not depend on \(h\).
\end{enumerate}
Let \(\hatt E\) and \(\hatt F\) be the principal-part polynomials of
\(E(f)\) and \(F(f)\). By (H2) they depend only on \(P\). Put
\begin{equation}\label{eq:phi-formal}
 \varphi(z;h)=\hatt E(D)\,F(f)-\hatt F(D)\,E(f) .
\end{equation}
By Lemma~\ref{lem:bracket} this is a power series in \(z\), and its
coefficients are polynomials in \(h\) and in the coefficients of \(P\).

\begin{lemma}\label{lem:bracket-wave}
Assume \textup{(H1)} and \textup{(H2)}, and let \(v\) be a wave. Then
the function
\[
 \Phi=\hatt E(D)\,F(u)-\hatt F(D)\,E(u),\qquad u=v/c,
\]
is entire, and at each pole \(z_i\) of \(v\),
\(\Phi(z_i+z)=\omega_i^{\,e}\,\varphi(z;h_i)\), with \(\omega_i,h_i\) as
in \eqref{eq:pole-germs}.
\end{lemma}

\begin{proof}
Near \(z_i\) we have \(u(z_i+z)=\omega_if(z;h_i)\). By (H1),
\(E(u)(z_i+z)=\omega_i^eE(f(\cdot\,;h_i))\) and similarly for \(F\).
Hence \(\Phi(z_i+z)=\omega_i^e\,\varphi(z;h_i)\), which is holomorphic.
Away from the poles, \(\Phi\) is holomorphic because \(u\) is.
\end{proof}

By Lemma~\ref{lem:bracket-wave} the values of \(\Phi\) at the poles are
\(\omega_i^e\,\varphi(0;h_i)\). If the constant term \(\varphi(0;h)\)
does not depend on \(h\), these values lie in a finite set, and
Lemma~\ref{lem:finite-values} makes \(\Phi\) constant. Then the coefficient of every positive power of \(z\) in
\(\varphi(z;h_i)\) vanishes. These are the equations for \(h_i\). In the example, where
\(\Phi=-\Psi'\), the constant term \(\varphi(0;h)\) is zero.

\subsection{The choice of the two expressions}\label{sec:choice}

It remains to choose \(E\) and \(F\). What we need from a pair is a
positive power of \(z\) whose coefficient in \(\varphi(z;h)\) contains
\(h\) with a nonzero numerical factor: on a wave this coefficient
vanishes, and this is an equation for \(h\). Both pairs
used below are built from the quadratic expression
\begin{equation}\label{eq:seed}
 H(u)=uu''-\frac{q+1}{q}(u')^2 ,
\end{equation}
which vanishes for \(u=z^{-q}\); its partner is a power of \(u\) with
the same behaviour under phase changes. For \(n\ge4\) the pair
\((u^2,H)\) gives an equation of degree one
(Section~\ref{sec:affine}). For \(n\le3\) that computation no longer
applies, and we use the pair \((u,H_n)\) with \(H_n=u^{n-2}H\). For this
pair the coefficients of \(\varphi\) below \(z^{J-1}\) have degree at
most one in \(h\), and their terms in \(h\) can all vanish, as they do
for \(D^4-s\); but the coefficient of \(z^{J-1}\) always contains
\(h^2\), and this leads to an equation of degree two
(Section~\ref{sec:quadratic}). The terms linear in \(h\) return in
Section~\ref{sec:not-even}, where they detect the odd powers in \(P\).

\begin{lemma}[The expression \(H\)]\label{lem:seed}
For \(u=z^{-q}(1+\varepsilon z^j)\),
\[
 H(u)=\varepsilon\,j(j+1)\,z^{j-2q-2}+O(\varepsilon^2).
\]
In particular \(H(z^{-q})=0\), so \(H(f)\) has a pole of order at most
\(2q+1\). Moreover
\[
 [h]\,H(f)=J(J+1)\,z^{p-2}\bigl(1+O(z)\bigr),\qquad
 [h^k]\,H(f)=O\bigl(z^{kJ-2q-2}\bigr)\quad(k\ge2).
\]
\end{lemma}

\begin{proof}
Put \(y=u^{-1/q}\), so that \(u=y^{-q}\). Then
\(u'=-q\,y^{-q-1}y'\) and
\(u''=q(q+1)\,y^{-q-2}(y')^2-q\,y^{-q-1}y''\), and therefore
\begin{equation}\label{eq:seed-y}
 H(u)=-q\,y^{-2q-1}\,y'' .
\end{equation}
For \(u=z^{-q}(1+\varepsilon z^j)\) we have
\(y=z-\tfrac{\varepsilon}{q}z^{j+1}+O(\varepsilon^2)\) and
\(y''=-\tfrac{\varepsilon}{q}\,j(j+1)z^{j-1}+O(\varepsilon^2)\), which
gives the first formula. The formulas for \(H(f)\) follow from it, with
\(j=J\), and from Lemma~\ref{lem:degree}, because \(H\) is quadratic
with two derivatives in total and \(J-2q-2=p-2\).
\end{proof}

\subsection{\texorpdfstring{The case \(n\ge4\): one value}{The case n >= 4: one value}}\label{sec:affine}

The pair \((u^2,H)\) requires \(p>2q+1\), that is, \((n-3)q>1\). For
even \(p\) this is the case exactly when \(n\ge4\): the inequality
holds for \(n\ge5\) and for \(n=4\), \(q\ge2\); it does not hold for
\(n\le3\); and for \(n=4\), \(q=1\) the degree \(p=3\) is odd.

\begin{proposition}[One value]\label{prop:affine}
Let \(p\) be even and \(n\ge4\), so that \(p>2q+1\).
Take \(E=u^2\) and \(F=H\). Then
\[
 [z^{p-2q-1}]\,\varphi(z;h)=\Delta_1\,(h-h_*),\qquad
 \Delta_1=-\frac{J(J+1)\,(p-2)!}{(2q-1)!\,(p-2q-1)!}\ne0,
\]
for a number \(h_*\) depending only on \(P\), and \(h_*\) is the free
coefficient at every pole of every wave. In particular, every wave has
one pole modulo \(\Gamma\), and any two waves coincide after
translation and phase change.
\end{proposition}

\begin{proof}
\emph{Hypotheses.} Both \(u^2\) and \(H\) are multiplied by
\(\omega^2\) under \(u\mapsto\omega u\). By Lemma~\ref{lem:degree} the
part of \(f^2\) depending on \(h\) is \(O(z^{J-2q})=O(z^p)\), and by
Lemma~\ref{lem:seed} that of \(H(f)\) is \(O(z^{p-2})\). Both are
holomorphic, so (H2) holds. The polynomial \(\hatt E\) has degree
\(2q-1\) and leading coefficient \(-1/(2q-1)!\), and \(\hatt F\) has
degree at most \(2q\).

\emph{The dependence on \(h\).} In
\([h]\varphi=\hatt E(D)\,[h]H(f)-\hatt F(D)\,[h]f^2\), the term of
highest order in \(\hatt E(D)\), namely \(-D^{2q-1}/(2q-1)!\), applied to
\(J(J+1)z^{p-2}\), gives \(\Delta_1z^{p-2q-1}\). All other terms of
\(\hatt E(D)[h]H(f)\) contain higher powers of \(z\). Since
\([h]f^2=2z^p(1+O(z))\) and \(\hatt F\) has degree at most \(2q\), the
second term is \(O(z^{p-2q})\). Hence
\[
 [h]\,\varphi=\Delta_1z^{p-2q-1}+O(z^{p-2q}).
\]
For \(k\ge2\), Lemmas~\ref{lem:degree} and~\ref{lem:seed} give
\([h^k]\varphi=O(z^{kJ-4q-1})\), and \(kJ-4q-1\ge2p-1>p-2q-1\). So the
coefficient of \(z^{p-2q-1}\) in \(\varphi\) is affine in \(h\) with
slope \(\Delta_1\ne0\). Because \(p-2q-1>0\), the constant term
\(\varphi(0)\) of \(\varphi\) does not depend on \(h\).

\emph{Conclusion.} By Lemma~\ref{lem:bracket-wave}, \(\Phi\) is
entire and its values at the poles lie in the finite set
\(\{\omega^2\varphi(0):\omega\in\mu_m\}\). By
Lemma~\ref{lem:finite-values}, \(\Phi\) is constant. The coefficient of
\(z^{p-2q-1}\) in \(\Phi(z_i+z)\), which is
\(\omega_i^2\Delta_1(h_i-h_*)\), therefore vanishes, and \(h_i=h_*\).
The last assertion follows from Corollary~\ref{cor:one-germ}.
\end{proof}

\subsection{\texorpdfstring{The cases \(n=2\) and \(n=3\): at most two values}{The cases n = 2 and n = 3: at most two values}}\label{sec:quadratic}

For \(n\le3\) we have \(p\le2q\); then \(D^{2q-1}\) annihilates
\(z^{p-2}\), and the term \(\Delta_1\) disappears. We take instead
\(E=u\) and \(F=H_n\), where
\[
 H_n(u)=u^{n-2}H(u)=u^{n-1}u''-\frac{q+1}{q}\,u^{n-2}(u')^2 ,
\]
and obtain an equation of degree two
(Proposition~\ref{prop:quadratic}). The proof uses an elementary
inequality for \(n=2\), which we state first.

\begin{lemma}\label{lem:kappa}
For every integer \(q\ge4\), the number \(X=(4q+1)_q/(q)_q\) satisfies
\(X>6q+3\).
\end{lemma}

\begin{proof}
Each of the \(q\) factors \((4q+1+i)/(q+i)\), \(0\le i<q\), of \(X\) is
at least \(5q/(2q-1)>5/2\), so \(X>(5/2)^q\). Moreover
\((5/2)^q>6q+3\) for \(q\ge4\): this holds for \(q=4\), and
\((5/2)^{q+1}-(6q+9)>\tfrac52(6q+3)-(6q+9)=9q-\tfrac32>0\).
\end{proof}

\begin{proposition}[A quadratic equation]\label{prop:quadratic}
Let \(p\ge4\) be even and \(q\ge2\); this covers every even \(p\ge4\)
with \(n=2\) or \(n=3\). There is a polynomial \(Q_P\) of degree two,
depending only on \(P\), such that \(Q_P(h)=0\) for the free
coefficient \(h\) at every pole of every wave.
\end{proposition}

\begin{proof}
Both \(u\) and \(H_n\) are multiplied by \(\omega\) under
\(u\mapsto\omega u\), since \(\omega^n=\omega\) for \(\omega\in\mu_m\);
so (H1) holds with \(e=1\). Let \(a_*=(-1)^{q-1}/(q-1)!\) be the
leading coefficient of \(\hatt E\). Write
\(f=f_0+hf_1+h^2f_2+O(h^3)\). By Lemma~\ref{lem:degree},
\[
 f_0=z^{-q}\bigl(1+O(z)\bigr),\qquad
 f_1=z^{nq}\bigl(1+O(z)\bigr),\qquad
 f_2=\psi\,z^{(2n+1)q}\bigl(1+O(z)\bigr),
\]
where \(\psi\) is the coefficient of \(h^2\) in \(b_{2J}\). By the
same lemma a monomial of \(b_{2J}\) that contains \(h^2\) has weight
\(2J\), so it is a multiple of \(h^2\) alone. In the recurrence at the
index \(2J\) such a term arises only from the \(\binom n2\) products
\(b_Jb_Jb_0^{\,n-2}\) in \(u^n\). So
\begin{equation}\label{eq:psi}
 \psi:=[h^2]\,b_{2J}=\binom n2\frac{K}{I(2J)}=\frac{\binom n2}{X-n},
 \qquad X=\frac{\bigl((n+2)q+1\bigr)_p}{(q)_p};
\end{equation}
here we used that the product \((2J-q)(2J-q-1)\cdots(2J-q-p+1)\) in
\(I(2J)\) equals \(((n+2)q+1)_p\). Since
\(((n+2)q+1)_p>(q+1)_p=n(q)_p\), we have \(X>n\) and \(\psi>0\).

The proof has five parts. The first three show that the terms of
\(\varphi\) of degree at least two in \(h\) begin with
\(\Delta_2h^2z^{J-1}\), where \(\Delta_2\ne0\); the last two derive the
quadratic equation from this.

\emph{Hypothesis \textup{(H2)}.} The part of \(E(f)=f\)
depending on \(h\) is \(O(z^{nq})\), hence holomorphic. The expression
\(H_n\) is a product of \(n\) factors with two derivatives in total; by
Lemma~\ref{lem:degree}, its part of degree \(k\ge1\) in \(h\) is
\(O(z^{kJ-nq-2})\), and for \(k=1\) the exponent is \(q-2\ge0\). So
(H2) holds. The polynomial \(\hatt E\) has degree \(q-1\), and
\(\hatt F\) has degree at most \(nq\) because \(H_n(f)\) has a pole of
order at most \((n-2)q+2q+1\).

\emph{The terms of \(H_n(f)\) in \(h\) and \(h^2\).} Let
\(\Bil\) be the symmetric bilinear form with \(\Bil(g,g)=H(g)\). Then
\([h]H(f)=2\Bil(f_0,f_1)\) and \([h^2]H(f)=2\Bil(f_0,f_2)+H(f_1)\), and
Lemma~\ref{lem:seed} gives
\[
 2\Bil(f_0,f_1)=J(J+1)\,z^{p-2}\bigl(1+O(z)\bigr),\qquad
 2\Bil(f_0,f_2)=2J(2J+1)\,\psi\,z^{2nq-2}\bigl(1+O(z)\bigr),
\]
while
\(H(f_1)=\bigl(nq(nq-1)-\tfrac{q+1}{q}(nq)^2\bigr)z^{2nq-2}(1+O(z))
=-nJ\,z^{2nq-2}(1+O(z))\). Recall that \(H(f_0)=O(z^{-2q-1})\). For the
term linear in \(h\),
\[
 [h]\,H_n(f)=f_0^{\,n-2}\cdot2\Bil(f_0,f_1)+(n-2)f_0^{\,n-3}f_1\,H(f_0)
 =J(J+1)\,z^{q-2}\bigl(1+O(z)\bigr),
\]
because the second term is \(O(z^{q-1})\). For the quadratic term,
\begin{align*}
 [h^2]\,H_n(f)={}&f_0^{\,n-2}\bigl(2\Bil(f_0,f_2)+H(f_1)\bigr)
 +(n-2)f_0^{\,n-3}f_1\cdot2\Bil(f_0,f_1)\\
 &+\Bigl((n-2)f_0^{\,n-3}f_2+\tbinom{n-2}{2}f_0^{\,n-4}f_1^{\,2}\Bigr)H(f_0).
\end{align*}
The first two terms begin with the power \(z^{(n+2)q-2}\), with the
coefficients \(2J(2J+1)\psi-nJ\) and \((n-2)J(J+1)\). The last term is
\(O(z^{(n+2)q-1})\), again because the leading pole of \(H(f_0)\)
cancels. Altogether,
\begin{equation}\label{eq:kappa}
 [h^2]\,H_n(f)=\kappa\,z^{(n+2)q-2}\bigl(1+O(z)\bigr),\qquad
 \kappa=J\bigl(2(2J+1)\psi+(n-2)J-2\bigr).
\end{equation}
For \(n\ge3\) we have \(\kappa>0\), since \(\psi>0\) and
\((n-2)J-2>0\). For \(n=2\) the degree \(p=q\) is even and at least
\(4\), \(J=3q\) and \(\psi=1/(X-2)\) with \(X=(4q+1)_q/(q)_q\), so that
\[
 \kappa=6q\Bigl(\frac{6q+1}{X-2}-1\Bigr)<0
\]
by Lemma~\ref{lem:kappa}. In all cases \(\kappa\ne0\).

\emph{The terms of \(\varphi\) of degree at least two in \(h\).} Only \(a_*D^{q-1}H_n(f)\) contributes to the lowest
\(h^2\)-term of \(\varphi\):
\begin{equation}\label{eq:Delta2}
 [h^2]\,\varphi=\Delta_2z^{J-1}+O(z^J),\qquad
 \Delta_2=a_*\,\kappa\,\frac{((n+2)q-2)!}{(J-1)!}\ne0.
\end{equation}
Indeed, \(D^{q-1}z^{(n+2)q-2}=\frac{((n+2)q-2)!}{(J-1)!}\,z^{J-1}\),
lower derivatives of \(H_n(f)\) give higher powers of \(z\), and since
\([h^2]f=O(z^{2J-q})\) and \(\hatt F\) has degree at most \(nq\), the
term \(\hatt F(D)[h^2]f\) is \(O(z^{2J-q-nq})=O(z^J)\). In the same way,
\([h^k]\varphi=O(z^{(k-1)J-1})\) for \(k\ge3\). So the terms of
\(\varphi\) of degree at least two in \(h\) start at \(z^{J-1}\), and
\begin{equation}\label{eq:slopes}
 [z^i]\,\varphi=A_ih+B_i\qquad(0\le i\le J-2),
\end{equation}
where \(A_i\) and \(B_i\) are polynomials in the coefficients
\(a_0,\dots,a_{p-1}\) of \(P\). In summary, the coefficients of \(h^k\) in
\(f\), \(H_n(f)\) and \(\varphi\) begin as follows; an entry without
\(O\) is a leading term, exact up to a factor \(1+O(z)\), and
\(\varphi\) is a power series by Lemma~\ref{lem:bracket}:
\begin{equation}\label{eq:leading-terms}
\begin{array}{c|ccc}
  k & f & H_n(f) & \varphi\\
 \hline
 0 & z^{-q} & O(z^{-nq-1}) & O(1)\\
 1 & z^{nq} & J(J+1)\,z^{q-2} & O(1)\\
 2 & \psi\,z^{(2n+1)q} & \kappa\,z^{(n+2)q-2} & \Delta_2\,z^{J-1}\\
 k\ge3 & O(z^{kJ-q}) & O(z^{kJ-nq-2}) & O(z^{(k-1)J-1})
\end{array}
\end{equation}

\emph{A function whose values at the poles do not depend on \(h\).} Let \(\Phi\) be the entire function of
Lemma~\ref{lem:bracket-wave}. At a pole \(z_i\) we have
\(\Phi(z_i)=\omega_i(A_0h_i+B_0)\) and, since \(J-2\ge1\),
\(\Phi'(z_i)=\omega_i(A_1h_i+B_1)\). Put
\[
 T=\Phi\ \text{ if } A_0=0,\qquad
 T=A_0\,\Phi'-A_1\,\Phi\ \text{ if } A_0\ne0.
\]
Then \(T\) is entire, and its value at a pole \(z_i\) is
\(\omega_iB_0\) in the first case and \(\omega_i(A_0B_1-A_1B_0)\) in the
second. These values lie in a finite set, so \(T\) is constant by
Lemma~\ref{lem:finite-values}. (By Lemma~\ref{lem:slopes}(b) below, the first case is \(a_{p-1}=0\);
this is not used in the proof.)

\emph{The quadratic equation.} If \(A_0=0\), the coefficient of
\(z^{J-1}\) in \(T(z_i+z)\) is \(\omega_i\) times a polynomial in
\(h_i\) of degree two with leading coefficient \(\Delta_2\). If
\(A_0\ne0\), the coefficient of \(z^{J-2}\) in \(T(z_i+z)\) is
\(\omega_i\) times a polynomial in \(h_i\) of degree two with leading
coefficient \(A_0(J-1)\Delta_2\): the derivative \(\Phi'\) contributes
\((J-1)\Delta_2h_i^2\) at this power, while \(\Phi\) has no
\(h^2\)-term below \(z^{J-1}\). In both cases the power of \(z\) is
positive, so the coefficient vanishes because \(T\) is constant.
Dividing by \(\omega_i\) gives the polynomial \(Q_P\).
\end{proof}


In Section~\ref{sec:selector-example}, \(A_0=0\),
\(T=\Phi=-\Psi'\), and \(Q_P\) is the coefficient of \(z^7\) in
\(-\Psi'\). Proposition~\ref{prop:quadratic} does not extend to
\(p=2\), where one equation can have infinitely many waves
(Appendix~\ref{sec:low-order}): at \((n,q)=(2,2)\) one finds \(X=15\)
and \(\kappa=0\).

% ======================================================================
\section{Proof of Theorem A and of the bounds in Theorem B}\label{sec:proof}

This section does Step~5 of Section~\ref{sec:example} in general and
proves Theorem~\ref{thm:A} and the bounds in Theorem~\ref{thm:B}.

\subsection{Two values in one wave}\label{sec:two-germs}

For \(n=2,3\) a wave might carry
the two roots of \(Q_P\) at different poles. A count of zeros and poles
shows that this does not happen.

\begin{lemma}[Two values, one orbit]\label{lem:two-germs}
Let \(p\) be even, and let \(v\) be a wave whose free coefficient takes
at most two values at its poles. Then \(v\) has exactly one pole modulo
\(\Gamma\).
\end{lemma}

\begin{proof}
The free coefficient determines the normalized series, so the series of
\(v\) fall into at most two classes. By Lemma~\ref{lem:germs}(a) its
poles lie in at most two cosets of \(\Gamma\), and by part~(c) it lies
in \(\Wclass\).
Suppose there are two, containing poles \(z_1\) and \(z_2\), and write
\(v(z_i+z)=c\,\omega_if(z;h_i)\) as in \eqref{eq:pole-germs}. Put
\[
 a=z_2-z_1,\qquad \omega=\frac{\omega_2}{\omega_1}\in\mu_m,\qquad
 W(x)=v(z_1+a+x)-\omega\,v(z_1+x).
\]
We show that \(W\equiv0\). Then \(a\in\Gamma\), which contradicts the
choice of \(z_1\) and \(z_2\). Suppose \(W\not\equiv0\).

\emph{Zeros and poles of \(W\).} By Corollary~\ref{cor:local}(b), \(W\)
has a zero of order at least \(nq\) at \(x=0\). It has the same
character as \(v\), so the same holds at every point of \(\Gamma\). The
poles of \(v\) form the set \((z_1+\Gamma)\cup(z_1+a+\Gamma)\), so the
poles of \(W\) lie in \(\Gamma\cup(a+\Gamma)\cup(-a+\Gamma)\). On
\(\Gamma\) the function \(W\) vanishes, so its poles lie in the two
cosets \(\pm a+\Gamma\), and they have order at most \(q\).

\emph{The compact surface.} Let \(\cyl\) be the Riemann sphere
\(\CC\cup\{\infty\}\) if \(\Gamma=\{0\}\); the cylinder \(\CC/\Gamma\)
with its two ends added, again a sphere, if \(\Gamma\) has rank one;
and the torus \(\CC/\Gamma\) if \(\Gamma\) is a lattice. The function
\(W^m\) is invariant under \(\Gamma\), so it is a meromorphic function
on \(\CC/\Gamma\). By Lemma~\ref{lem:germs}(c), \(v\) and hence \(W^m\)
are bounded near the added points, so \(W^m\) extends to a meromorphic
function on \(\cyl\). The quotient map is locally biholomorphic, so it
preserves orders of zeros and poles. A nonzero meromorphic function on
a compact Riemann surface has as many zeros as poles, counted with
multiplicity.

\emph{The case \(n\ge3\).} On \(\cyl\) the function \(W^m\) has poles of
total order at most \(2qm\), and a zero of order at least \(nqm\). This
is a contradiction.

\emph{The case \(n=2\).} Now \(m=1\) and \(\omega=1\), the leading
coefficient \(c\) is the same at all poles, and \(q=p\) is even. If
\(2a\in\Gamma\), the two cosets \(\pm a+\Gamma\) coincide; then \(W\)
has poles of total order at most \(q\) and a zero of order at least
\(2q\), a contradiction. If \(2a\notin\Gamma\), this count gives no
contradiction: \(W\) has poles of total order at most \(2q\) and a zero
of order at least \(2q\). One more zero comes from the odd part of
\(W\). Consider
\[
 \Omega(x)=W(x)-W(-x)
 =v(z_1+a+x)-v(z_1+x)-v(z_1+a-x)+v(z_1-x).
\]
The reflection \(x\mapsto-x\) acts on \(\cyl\), so \(\Omega\) is
meromorphic there, with poles only at the two points \(\pm a\), of
order at most \(q\). It is odd and \(O(x^{2q})\) at \(0\), hence
\(O(x^{2q+1})\). At \(x=-a\) the four terms are evaluated at \(z_1\),
\(z_1-a\), \(z_1+2a\) and \(z_1+a\). The points \(z_1-a\) and
\(z_1+2a\) are not poles of \(v\), since neither \(a\) nor \(2a\) lies
in \(\Gamma\). So near \(x=-a\) only the term \(v(z_1+a+x)\) of
\(W(x)\) has a pole, with leading coefficient \(c\), and only the term
\(-v(z_1-x)\) of \(W(-x)\) has a pole, with leading coefficient \(-c\)
because \(q\) is even. Thus \(\Omega(x)=2c\,(x+a)^{-q}+\cdots\) is not
identically zero, and on \(\cyl\) it has at least \(2q+1\) zeros and at
most \(2q\) poles. This is impossible. So \(W\equiv0\) in all cases.
\end{proof}

\subsection{Proof of Theorem A, and the bound of two waves}

Let \(v\) be a wave. By Proposition~\ref{prop:existence} it has a pole,
so \(q=p/m\) is a positive integer. For odd \(p\) the theorem was
proved in Section~\ref{sec:odd}. Let \(p\) be even. If \(n\ge4\),
Proposition~\ref{prop:affine} shows that all poles of all waves carry
the same free coefficient, and Corollary~\ref{cor:one-germ}(a) applies.
If \(n\in\{2,3\}\), then \(p\ge4\), because \(p\ge3\) is even, and
\(q\ge2\); by Proposition~\ref{prop:quadratic} the free coefficients at
the poles of all waves are roots of \(Q_P\), and
Lemma~\ref{lem:two-germs} applies.
In each case \(v\) has exactly one pole modulo \(\Gamma\), and
Proposition~\ref{prop:spaces} gives its form. This proves
Theorem~\ref{thm:A}.

By Corollary~\ref{cor:one-germ}(b), two waves that share a value of the
free coefficient coincide after translation and phase change. So an
equation has at most as many waves as there are values of the free
coefficient: at most two, and only one if \(n\ge4\). This proves the
first assertion of Theorem~\ref{thm:B}.

\subsection{Operators that are not even}\label{sec:not-even}

For \(n=2,3\) the quadratic equation leaves two possible values of the
free coefficient. The terms of \(\varphi\) that are linear in \(h\)
decide between them whenever \(P\) is not even. We keep the pair
\(E=u\), \(F=H_n\) and the assumptions \(p\ge4\) even and \(q\ge2\), and
recall from \eqref{eq:slopes} that \([z^i]\varphi=A_ih+B_i\) for
\(0\le i\le J-2\), and that \(a_{p-k}\) has weight \(k\).

\begin{lemma}[The terms linear in \(h\)]\label{lem:slopes}
Let \(p\ge4\) be even and \(q\ge2\).
\begin{enumerate}[label=\textup{(\alph*)},leftmargin=2.2em]
\item \(A_i\) is a weighted-homogeneous polynomial of weight \(i+1\).
\item \(A_0=c_0\,a_{p-1}\) and \(A_1=c_1\,a_{p-1}^2\), where
\[
 c_0=(-1)^q\,\frac{2(n+1)(nq+1)(J+1)}{(n-1)(J-1)},\qquad
 c_1=-\frac{nq+2}{(n-1)(J-1)}\,c_0 .
\]
\item Let \(n\in\{2,3\}\), let \(j\) be odd with \(3\le j\le p-1\), and
suppose that \(a_{p-i}=0\) for all odd \(i<j\). Then
\(A_{j-1}=c_{j-1}\,a_{p-j}\) with a number \(c_{j-1}\ne0\) depending
only on \(n\), \(q\) and \(j\).
\end{enumerate}
\end{lemma}

\begin{proof}
(a) By the proof of Lemma~\ref{lem:degree}, \(\lambda^qf(\lambda z;h)\)
is the series \eqref{eq:family} of the operator \(P_\lambda\) with
free coefficient \(\lambda^Jh\). The expressions \(E=u\) and \(F=H_n\)
are homogeneous: \(H_n\) has \(n\) factors and two derivatives, so
\(H_n\bigl(\lambda^qf(\lambda\,\cdot)\bigr)(z)
=\lambda^{nq+2}H_n(f)(\lambda z)\). If a Laurent series \(\alpha\) has the
principal-part polynomial \(\hatt\alpha(\rho)\), then
\(\lambda^N\alpha(\lambda z)\) has
\(\lambda^{N-1}\hatt\alpha(\rho/\lambda)\), for every integer \(N\).
Inserting this into \eqref{eq:phi-formal} shows that the series
\(\varphi\) of \(P_\lambda\), taken at the free coefficient
\(\lambda^Jh\), equals \(\lambda^{J+1}\varphi(\lambda z;h)\). Comparing
the coefficients of \(h\,z^i\) gives
\(A_i(P_\lambda)=\lambda^{i+1}A_i(P)\), and the coefficients of
\(P_\lambda\) are \(\lambda^ka_{p-k}\).

(b), (c) By (a), \(A_0\) is a multiple of \(a_{p-1}\), and \(A_1\) is a
combination of \(a_{p-1}^2\) and \(a_{p-2}\). Under the hypothesis of
(c), \(A_{j-1}\) has odd weight \(j\), so each of its monomials contains
a coefficient \(a_{p-i}\) with \(i\) odd; those with \(i<j\) vanish,
and the only monomial of weight \(j\) containing \(a_{p-i}\) with
\(i\ge j\) is \(a_{p-j}\) itself. It remains to find the numerical
coefficients. They depend only on \(n\), \(q\) and \(j\). Since
\([z^i]\varphi=A_ih+B_i\), the number \(A_i\) is the coefficient of
\(z^i\) in \(\partial\varphi/\partial h\) at every value of \(h\). So
the coefficients can be computed from any convenient operator and any
value of \(h\). This is done in
Appendix~\ref{app:slopes} by three explicit computations with Laurent
series. First, \(A_1=0\) for every even operator, so \(A_1\) contains
no term \(a_{p-2}\). Second, \(c_0\) and \(c_1\) are read off from an
operator that has an explicit wave; at the free coefficient of this
wave the derivative of \(f(z;h)\) with respect to \(h\) is also
explicit. Third, \(c_{j-1}\) is found from the operator
\(D^p+a\,D^{p-j}\) to first order in \(a\), and an inequality between
binomial coefficients shows that it is not zero.
\end{proof}

\begin{proposition}[One wave unless \(P\) is even]\label{prop:not-even}
Let \(p\ge4\) be even and \(n\in\{2,3\}\).
\begin{enumerate}[label=\textup{(\alph*)},leftmargin=2.2em]
\item If \(A_i\ne0\) for some \(i\) with \(1\le i\le J-2\), then the
free coefficient equals \(-B_i/A_i\) at all poles of all waves, and
\eqref{eq:main} has at most one wave up to translation and phase
change.
\item If \(P\) is not even, then \(A_i\ne0\) for some \(i\) with
\(1\le i\le p-2\).
\end{enumerate}
\end{proposition}

\begin{proof}
(a) Let \(v\) be a wave, and let \(\Phi\) be the entire function of
Lemma~\ref{lem:bracket-wave} for \(E=u\), \(F=H_n\). By
Proposition~\ref{prop:quadratic} the free coefficients \(h_i\) at the
poles of \(v\) are roots of \(Q_P\). So the values
\(\omega_i(A_0h_i+B_0)\) of \(\Phi\) at the poles lie in a finite set,
and \(\Phi\) is constant by Lemma~\ref{lem:finite-values}. Hence the
coefficients of all positive powers of \(z\) in \(\varphi(z;h)\) vanish
for the free coefficient \(h\) at any pole of any wave; by
\eqref{eq:slopes},
\[
 A_ih+B_i=0\qquad(1\le i\le J-2).
\]
If \(A_i\ne0\) for one of these \(i\), then \(h=-B_i/A_i\). So all
poles of all waves carry the same normalized series, and
Corollary~\ref{cor:one-germ} applies.

(b) Since \(P\) is not even and \(p\) is even, there is a least odd
\(j\) with \(a_{p-j}\ne0\), and \(j\le p-1\). If \(j=1\), then
\(A_1=c_1a_{p-1}^2\ne0\). If \(j\ge3\), then
\(A_{j-1}=c_{j-1}a_{p-j}\ne0\), and \(2\le j-1\le p-2\).
\end{proof}

So an equation has at most one wave if \(n\in\{2,3\}\) and \(P\) is not
even. Together with Section~\ref{sec:odd} and
Proposition~\ref{prop:affine} this proves the second assertion of
Theorem~\ref{thm:B}; Table~\ref{tab:cases} collects the cases. The last
assertion, that \(D^4-s\) with \(s\ne0\) has exactly two waves, is
proved in
Section~\ref{sec:two-waves} for \(n=3\) and in Section~\ref{sec:n2} for
\(n=2\).

Proposition~\ref{prop:not-even}(a) can also apply to an even operator
when \(p\ge6\). For
\(n=3\) and \(P(D)=D^6+a_0\) with \(a_0\ne0\) one finds
\(A_5=-\tfrac{13}{888}\,a_0\), so the equation
\(v^{(6)}+a_0v+v^3/3=0\) has only one wave,
\(v=-\tfrac c2\,\wp'(z;0,-a_0/4320)\).

\begin{table}[ht]
\centering
\small
\begin{tabular}{@{}lll@{}}
\toprule
case & values of \(h\) on all waves of one equation & proved in\\
\midrule
\(p\) odd & no free coefficient & Lemma~\ref{lem:resonance}\\
\(p\) even, \(n\ge4\) & at most one & Proposition~\ref{prop:affine}\\
\(p\) even, \(n=2,3\), \(P\) not even & at most one &
 Propositions~\ref{prop:quadratic} and~\ref{prop:not-even}\\
\(p\) even, \(n=2,3\), \(P\) even & at most two &
 Proposition~\ref{prop:quadratic}\\
\bottomrule
\end{tabular}
\caption{The free coefficient \(h\) in the proof of
Theorems~\ref{thm:A} and~\ref{thm:B}. In each case a wave has one pole
modulo \(\Gamma\) (Corollary~\ref{cor:one-germ}(a) in the first two
rows, Lemma~\ref{lem:two-germs} in the last two) and therefore carries one
value of \(h\) (Lemma~\ref{lem:germs}(a)); each value belongs to at
most one wave
(Corollary~\ref{cor:one-germ}(b)).}
\label{tab:cases}
\end{table}

% ======================================================================
\section{The power of a wave}\label{sec:power}

This section does Step~7 of Section~\ref{sec:example} in general and
proves Theorem~\ref{thm:C}.

\subsection{The formula}\label{sec:power-formula}

Let \(p\ge3\) and let \(u=z^{-q}\sum_jb_jz^j\) be a normalized wave. The coefficients \(b_1,\dots,b_p\) precede the only
possible resonance \(J=p+2q\), so they are determined by \(P\), and the
recurrence defines them for every operator, whether or not a wave
exists. Define \(\ell_0=1,\ell_1,\dots,\ell_p\) by
\begin{equation}\label{eq:ell}
 \Bigl(\sum_{i=0}^{p}b_i(P)\,z^i\Bigr)^{m}
 =\sum_{j=0}^{p}\ell_j(P)\,z^j+O(z^{p+1}).
\end{equation}
Then every normalized wave satisfies
\[
 u^m=\sum_{j=0}^{p}\ell_j(P)\,z^{j-p}+O(z),
\]
so both the principal part and the constant term of \(u^m\) are fixed
by \(P\).

The power \(u^m\) has trivial character, because \(\chi\) takes values
in \(\mu_m\); so \(u^m\) is invariant under \(\Gamma\), with one pole
of order \(p\) modulo \(\Gamma\). We shall write it as a combination
of standard functions \(\Lambda_r\). For each group \(\Gamma\), and for
\(r\) in the range given below, there is a unique function
\(\Lambda_r\) with the following properties. It is invariant under
\(\Gamma\), has poles only on \(\Gamma\), and satisfies
\(\Lambda_r=z^{-r}+O(z)\) at \(0\). It is bounded at infinity if
\(\Gamma=\{0\}\), and at the two ends if \(\Gamma\) has rank one.
Explicitly:
\begin{itemize}
\item if \(\Gamma=\{0\}\): \(\Lambda_r=z^{-r}\), \(1\le r\le p\);
\item if \(\Gamma=(2\pi i/k)\ZZ\): with \(G=k\,e^{kz}/(e^{kz}-1)\),
\[
 \Lambda_r=\frac{(-1)^{r-1}}{(r-1)!}
 \Bigl(D^{r-1}G-[z^0]\,D^{r-1}G\Bigr),\qquad 1\le r\le p;
\]
\item if \(\Gamma\) is a lattice: with \(\wp=\wp(z;\Gamma)\),
\[
 \Lambda_r=\frac{(-1)^{r}}{(r-1)!}
 \Bigl(D^{r-2}\wp-[z^0]\,D^{r-2}\wp\Bigr),\qquad 2\le r\le p.
\]
\end{itemize}
Since \(D^jG=(-1)^jj!\,z^{-j-1}+O(1)\) and
\(D^j\wp=(-1)^j(j+1)!\,z^{-j-2}+O(1)\), these functions indeed satisfy
\(\Lambda_r=z^{-r}+O(z)\). Each \(\Lambda_r\) is unique, because the
difference of two candidates is a bounded function without poles on
\(\CC/\Gamma\) (or on the sphere), hence a constant, and its constant
term is zero. No elliptic function has a single simple pole, so there
is no \(\Lambda_1\) in the lattice case. Writing a function of
\(\Wclass\) through the derivatives of one function, here \(G\) or
\(\wp\), is Hermite's decomposition, a standard tool for constructing
meromorphic travelling waves
\citep{MusetteConte2003,DeminaKudryashov2011,DeminaKudryashovElliptic2011,DeminaKudryashovKawahara2010}.

\begin{theorem}[The formula for \(u^m\)]\label{thm:power}
Let \(p\ge3\) and let \(u\) be a normalized wave with group of twisted
periods \(\Gamma\). Then
\begin{equation}\label{eq:completion}
 u^m=\ell_p+\sum_{r=1}^{p}\ell_{p-r}\Lambda_r,
\end{equation}
with the functions \(\Lambda_r\) of the group \(\Gamma\), and:
\begin{enumerate}[label=\textup{(\roman*)},leftmargin=2.2em]
\item if \(\Gamma=\{0\}\), then \(\ell_p=0\);
\item if \(\Gamma=(2\pi i/k)\ZZ\), then the right-hand side is a monic
polynomial of degree \(p\) in \(G\); the residue \(\ell_{p-1}\) may be
nonzero (Corollary~\ref{cor:residue});
\item if \(\Gamma\) is a lattice, then \(\ell_{p-1}=0\), the term
\(r=1\) is absent, and the right-hand side is an elliptic function for
\(\Gamma\).
\end{enumerate}
In each case the right-hand side is the only function that is invariant
under \(\Gamma\), has poles only on \(\Gamma\), is bounded at infinity or
at the two ends when \(\Gamma\) has rank zero or one, and has the
principal part and constant term prescribed by \eqref{eq:ell}. In particular, \(u^m\) does not depend on the character
or on any further coefficient of the wave.
\end{theorem}

\begin{proof}
By Theorem~\ref{thm:A}, \(u^m\) is invariant under \(\Gamma\), has a
single pole of order \(p\) modulo \(\Gamma\), and is bounded at infinity
or at the two ends when \(\Gamma\) has rank zero or one.

(i) The difference \(u^m-\sum_r\ell_{p-r}z^{-r}\) has no poles and is
bounded, so it is a constant, equal to its constant term \(\ell_p\).
Both \(u^m\) and the sum tend to \(0\) at infinity by
Theorem~\ref{thm:A}(i), so \(\ell_p=0\).

(ii) By \eqref{eq:ell}, \(u^m\) and the right-hand side of
\eqref{eq:completion} have the same principal part and constant term
at \(0\). Their difference is invariant under \(\Gamma\), has no poles
and is bounded at both ends; so it is a constant, which is zero. Since
\(DG=G(k-G)\), the function \(D^{r-1}G\) is a polynomial of degree \(r\)
in \(G\) with leading coefficient \((-1)^{r-1}(r-1)!\); so \(\Lambda_p\)
is monic of degree \(p\) in \(G\), and \(\ell_0=1\).

(iii) The residue \(\ell_{p-1}\) of the elliptic function \(u^m\) at its
only pole in a period cell vanishes. As in (ii), the difference of
\(u^m\) and the right-hand side is then an elliptic function without
poles, hence a constant, which is zero.

Uniqueness, among the functions that are invariant under \(\Gamma\),
have poles only on \(\Gamma\) and are bounded at infinity or at the
ends, follows as for \(\Lambda_r\).
\end{proof}

All coefficients in these formulas are explicit. For instance,
\([z^0]D^jG=B_{j+1}(1)\,k^{j+1}/(j+1)\) with the Bernoulli polynomials
\(B_j\), and the constant terms of \(D^j\wp\) are polynomials in the
invariants of the lattice. The theorem describes \(u^m\) once \(\Gamma\)
is known; it does not determine the rate \(k\) or the lattice.

\begin{corollary}[Power residue]\label{cor:residue}
Let \(p\ge3\). The number \(\mathcal R(P)=\ell_{p-1}(P)\) is a polynomial in the
coefficients of \(P\), and for a normalized wave it equals
\(\Res_{z=0}u^m\). Every elliptic wave requires \(\mathcal R(P)=0\).
For a wave with \(\Gamma=(2\pi i/k)\ZZ\),
\begin{equation}\label{eq:end-values}
 u_+^{\,m}-u_-^{\,m}=k\,\mathcal R(P),
\end{equation}
where \(u_\pm\) are the limits of \(u\) as
\(\operatorname{Re}(kz)\to\pm\infty\).
\end{corollary}

\begin{proof}
The elliptic case is part (iii) of Theorem~\ref{thm:power}. If \(\Gamma\)
has rank one, then \(G\to k\) as \(\operatorname{Re}(kz)\to+\infty\)
and \(G\to0\) as \(\operatorname{Re}(kz)\to-\infty\). Since
\(DG=G(k-G)\), every derivative \(D^jG\), \(j\ge1\), vanishes at both
ends, so \(\Lambda_r\) takes the same value at both ends for \(r\ge2\),
while \(\Lambda_1=G-k/2\) increases by \(k\). Formula
\eqref{eq:completion} gives \eqref{eq:end-values}.
\end{proof}

For elliptic solutions, conditions obtained from the vanishing of a sum
of residues are a standard tool \citep{Hone2005,ConteNg2010}; the
corollary applies it to \(u^m\), which has one pole modulo \(\Gamma\).
Theorem~\ref{thm:power} and Corollary~\ref{cor:residue} are the precise
form of Theorem~\ref{thm:C}, which is thereby proved. Formula
\eqref{eq:completion} has further consequences.

\begin{corollary}\label{cor:power-consequences}
Let \(p\ge3\).
\begin{enumerate}[label=\textup{(\alph*)},leftmargin=2.2em]
\item A rational wave has \(\ell_{p-r}=0\) for \(1\le r<m\), besides
\(\ell_p=0\).
\item A twisted wave with \(\Gamma\) of rank one has
\(\mathcal R(P)=0\).
\item If \(n\ge3\) and \(\mathcal R(P)\ne0\), every wave is an
untwisted front: \(\Gamma\) has rank one, the character is trivial, and
the limits at the two ends are different.
\item Two waves of one equation that are not related by translation
and phase change never have the same twisted periods.
\end{enumerate}
\end{corollary}

\begin{proof}
(a) By Theorem~\ref{thm:power}(i), \(u^m=\sum_r\ell_{p-r}z^{-r}\), and
\(u^m=O(z^{-m})\) at infinity by Theorem~\ref{thm:A}(i).

(b) Such a wave tends to \(0\) at both ends
(Theorem~\ref{thm:A}(ii)), so \(\mathcal R(P)=0\) by
\eqref{eq:end-values}.

(c) Now \(m\ge2\) and \(\ell_{p-1}=\mathcal R(P)\ne0\). By (a) with
\(r=1\) no wave is rational, by Corollary~\ref{cor:residue} none is
elliptic, and by (b) none is twisted. So \(\Gamma\) has rank one and
the character is trivial, and \(u_+^m\ne u_-^m\) by
\eqref{eq:end-values}.

(d) If two normalized waves of the same equation have the same group
\(\Gamma\), then \eqref{eq:completion} gives them the same power
\(u^m\); their quotient is then a constant in \(\mu_m\), which is \(1\)
because both begin with \(z^{-q}\).
\end{proof}

By part~(d), a wave is determined, up to a phase change, by \(P\), the
position of one pole and the group \(\Gamma\). If \(\Gamma\) has rank
one and \(\mathcal R(P)\ne0\), the limits satisfy
\(P(0)u_\pm=Ku_\pm^{\,n}\) by \eqref{eq:normalized}, so
\(u_\pm^m\in\{0,P(0)/K\}\), and \eqref{eq:end-values} gives the rate
\(k=\pm P(0)/(K\mathcal R(P))\).

\subsection{The residue condition: the Kuramoto--Sivashinsky equation}\label{sec:ks}\label{ex:ks}

The Kuramoto--Sivashinsky equation with dispersion,
\(U_t+UU_x+\mu U_{xx}+\delta U_{xxx}+\nu U_{xxxx}=0\), \(\nu\ne0\),
also describes pressure waves in a liquid with gas bubbles
\citep{KudryashovSinelshchikovFourthOrder2010}. Its travelling waves
\(U=w(x-st)\) satisfy
\[
 \nu w'''+\delta w''+\mu w'-sw+\frac{w^2}{2}+C=0
\]
for a constant \(C\). Since \(n=2\), the shift \(w=v+w_0\) with
\(w_0^2/2-sw_0+C=0\) removes \(C\) and changes only the constant
term of the linear part: \(v\) satisfies \(P(D)v+v^2/2=0\) for the operator
\(P(D)=\nu D^3+\delta D^2+\mu D+a_0\) with \(a_0=w_0-s\), and \(q=3\).

Consider first \(\nu=1\), so that \(P\) is monic. Then \(c=120\), and
\(u=v/120\) satisfies \(P(D)u+60u^2=0\). For
\(u=z^{-3}+b_1z^{-2}+b_2z^{-1}+O(1)\), the coefficients of \(z^{-5}\)
and \(z^{-4}\) give
\[
 96b_1+12\delta=0,\qquad 114b_2+6\delta b_1-3\mu+60b_1^2=0,
\]
so that
\[
 u=z^{-3}-\frac{\delta}{8}z^{-2}+\frac{16\mu-\delta^2}{608}z^{-1}+O(1).
\]
Here \(m=1\) and \(\mathcal R(P)=\Res_{z=0}u=(16\mu-\delta^2)/608\).
For general \(\nu\), the function \(v/\nu\) solves the equation of the
monic operator \(P/\nu\). So \(v=120\nu u\), where \(u\) is a normalized
wave of \(P/\nu\), and
\(\mathcal R=(16\mu\nu-\delta^2)/(608\nu^2)\). By
Corollary~\ref{cor:residue}, every elliptic travelling wave requires
\[
 \delta^2=16\mu\nu,
\]
independently of \(s\) and \(C\)
\citep[Eq.~(3) and Theorem~1(i)]{Eremenko2006}.

The condition holds for \(\delta=4\) and \(\mu=\nu=1\). At the speed
\(s=0\), the functions
\[
 w=-60\,\wp'-60\,\wp-1,\qquad \wp=\wp\bigl(z;\tfrac1{12},g_3\bigr),
\]
solve \(w'''+4w''+w'+w^2/2+C=0\) with \(C=1080g_3-13\). For
\(27g_3^2\ne12^{-3}\) they are elliptic. So at a single speed there
are infinitely many elliptic travelling waves,
one for each integration constant other than \(-8\) and \(-18\). This
is why the bound of Theorem~\ref{thm:B} refers to a fixed integration
constant as well as to a fixed speed: the operator \(P\) depends on
both. For \(C=-8\) and \(C=-18\), that is, \(g_3=\pm1/216\), the
function \(\wp\) degenerates, and the formula gives a wave that is
rational in an exponential. By Theorem~\ref{thm:B} it is the only
wave, so for these two constants there is no elliptic wave. Thus the
condition \(\delta^2=16\mu\nu\) is necessary for an elliptic wave and
not sufficient. The elliptic travelling waves of the
Kuramoto--Sivashinsky equation with dispersion were found by
\citet{Kudryashov1990}, and by \citet{Eremenko2006} there are no others.

\subsection{\texorpdfstring{Two waves for \(n=2\)}{Two waves for n = 2}}\label{sec:n2}

The bound of two waves is also attained for \(n=2\), and
Theorem~\ref{thm:power} finds the waves. For
\(P(D)=D^4-s\) with \(s\ne0\) we have \(q=4\), \(m=1\) and
\(K=840\), and the normalized equation is \(u''''=s\,u+840\,u^2\). The
recurrence gives \(u=z^{-4}-\tfrac{s}{1680}+O(z^{4})\), so that
\(\ell_1=\ell_2=\ell_3=0\) and \(\ell_4=-s/1680\) in \eqref{eq:ell}. By
Theorem~\ref{thm:power}, every normalized wave is
\(u=\Lambda_4-s/1680\). To write \(\Lambda_4\) in the same way for
all three kinds of \(\Gamma\), we use the Weierstrass function
\(\wp=\wp(z;g_2,g_3)\) and its degenerations. If \(\Gamma\) is a
lattice, \(\wp\) is its Weierstrass function. If
\(\Gamma=(2\pi i/k)\ZZ\), put \(\wp=k^2/12-DG\) and
\((g_2,g_3)=(k^4/12,-k^6/216)\). If \(\Gamma=\{0\}\), put
\(\wp=z^{-2}\) and \((g_2,g_3)=(0,0)\). In all three cases
\(\wp''=6\wp^2-g_2/2\) and
\(\Lambda_4=\tfrac16\bigl(\wp''-[z^0]\wp''\bigr)=\wp^2-g_2/10\), so
that
\[
 u=\wp^2-\frac{g_2}{10}-\frac{s}{1680}.
\]
Substituting into the equation gives
\[
 g_3=0,\qquad g_2^2=\frac{s^2}{23184}.
\]
So \(g_3=0\) and \(g_2\ne0\), and \(\wp\) is not degenerate. Hence \(v''''-sv+v^2/2=0\) has exactly two waves for every
\(s\ne0\), both elliptic, one for each root
\(g_2=\pm s/(12\sqrt{161})\); they are different waves, because their
free coefficients \(h=g_2^3/10400\) are opposite and not zero. They are the
solutions (26) of \citet{DeminaKudryashovKawahara2010}, with \(g_2\) a
root of the quadratic equation (25) of that paper, at \(\beta=0\) and
\(C_1=0\) in its notation. With Section~\ref{sec:two-waves}, this
completes the proof of Theorem~\ref{thm:B}.

The same computation shows that the bound is attained for \(p=6\). For
\(P(D)=D^6+a_0\) with \(a_0\ne0\) it gives
\(u=\wp^3-\tfrac{3g_2}{20}\wp-\tfrac{3g_3}{28}+\tfrac{a_0}{665280}\)
with \(g_2=0\) and \(g_3^2=a_0^2/3643833600\). So
\(v^{(6)}+a_0v+v^2/2=0\) has exactly two waves, both elliptic, one for
each sign of \(g_3\); they are different waves, because their lattices
are different.
% ======================================================================
\section{Further consequences, and an example}\label{sec:consequences}\label{sec:examples}

\subsection{Equations without dissipation}\label{sec:even}

If \(v\) solves \eqref{eq:main}, then \(v(-z)\), multiplied by a number
whose \(m\)th power is \((-1)^p\), solves the same equation with
\(P(\rho)\) replaced by \((-1)^pP(-\rho)\). In particular, for even
\(P\) the \emph{reflection} \(v(z)\mapsto v(-z)\) is a symmetry of
\eqref{eq:main}, and every wave is symmetric:

\begin{proposition}[Reflection]\label{prop:reflection}
Let \(P\) be even. Every wave \(v\) has one pole modulo \(\Gamma\), its
character has order at most two, and at every pole \(z_0\)
\begin{equation}\label{eq:reflection}
 v(z_0+z)=(-1)^q\,v(z_0-z).
\end{equation}
\end{proposition}

\begin{proof}
The pole orbit comes from Theorem~\ref{thm:A} if \(p\ge3\) and from
Proposition~\ref{prop:low-order} if \(p=2\). Since \(P\) is even,
\(v(2z_0-z)\) is again a wave. If \(q\) is odd, then \(m=p/q\) is even,
so \(-1\in\mu_m\) and \(-v(2z_0-z)\) is a wave too. In both cases
\(\tilde v(z)=(-1)^qv(2z_0-z)\) is a wave with the same leading
coefficient as \(v\) at \(z_0\). Reflection multiplies the normalized
Laurent coefficient \(b_j\) by \((-1)^j\), so it preserves the free
coefficient \(b_J\), because \(J=p+2q\) is even. A normalized series is
determined by its free coefficient, so the series of \(v\) and
\(\tilde v\) at \(z_0\) coincide, and \(\tilde v=v\). This is
\eqref{eq:reflection}.

Place the pole at \(0\), so that \(v(-z)=(-1)^qv(z)\). For a twisted
period \(\gamma\), \(v(z+\gamma)=\chi(\gamma)v(z)\) and reflection give
\(v(z-\gamma)=\chi(\gamma)v(z)\). On the other hand
\(v(z-\gamma)=\chi(\gamma)^{-1}v(z)\). Hence \(\chi(\gamma)^2=1\).
\end{proof}

\begin{corollary}[Limits at the two ends]\label{cor:ends}
Let \(P\) be even and let \(v\) be a wave with \(\Gamma\) of rank one.
The limits \(v_\pm\) of \(v\) at its two ends are constant solutions of
\eqref{eq:main}, and \(v_+=(-1)^qv_-\). Hence a front, with
\(v_+\ne v_-\), requires odd \(q\); then \(n\) is odd, and \(v_\pm\) are
opposite nonzero constant solutions.
\end{corollary}

\begin{proof}
By Proposition~\ref{prop:reflection} and Lemma~\ref{lem:germs}(c),
\(v\) is a rational function of \(t=e^{2\pi iz/\tau}\) that is finite
at \(t=0\) and \(t=\infty\);
here \(\tau\) generates the group of ordinary periods. So
the limits exist, and all derivatives of \(v\) tend to \(0\) at the
ends, because \(D=(2\pi i/\tau)\,t\,d/dt\). Letting
\(\operatorname{Re}(kz)\to\pm\infty\) in \eqref{eq:main} shows that
\(v_\pm\) are constant solutions. Reflection in a pole interchanges the
two ends, so taking limits in \eqref{eq:reflection} gives
\(v_+=(-1)^qv_-\). If \(q\) is odd, \(n\) is odd as in the previous
proof, and a front has \(v_+=-v_-\ne v_-\), so \(v_\pm\ne0\).
\end{proof}

\subsection{A first-derivative term}\label{sec:drift}

If a first-derivative term is added to an even operator, the
reflection symmetry is lost. For \(n\ge3\) and \(p\ge4\) there are then
no waves at all.

\begin{corollary}[No waves with a first-derivative term]\label{cor:drift}
Let \(n\ge3\), let \(p\ge4\) be even, and let \(P(D)=P_e(D)+a_1D\),
where \(P_e\) is even and \(a_1\ne0\). Then every meromorphic solution
of \eqref{eq:main} is constant; the constant solutions are \(v=0\) and
the roots of \(v^m=-nP(0)\).
\end{corollary}

\begin{proof}
Suppose there is a wave; then \(p=mq\) with \(q\) an integer, and
\(K=(q)_p\). At a pole, in the normalization \eqref{eq:normalized},
induction on \eqref{eq:recurrence} gives \(b_j=0\) for odd \(j<p-1\): the even part \(P_e\) and the
nonlinearity preserve the parity of Laurent indices, and the term
\(a_1D\) first contributes at index \(p-1\). At that index
\(I(p-1)=-nK\), because \(0\le p-1-q<p\). Hence
\[
 b_{p-1}=-\frac{qa_1}{nK},\qquad
 \mathcal R(P)=\Res_{z=0} u^m=m\,b_{p-1}=-\frac{pa_1}{nK}\ne0 ;
\]
here \(\Res u^m=mb_{p-1}\) because \(b_{p-1}\) is the only coefficient
of odd index below \(p\). Since \(b_j=0\) for odd \(j<p-1\),
\begin{equation}\label{eq:drift-parity}
 v(z_0+z)-(-1)^qv(z_0-z)=O(z^{p-1-q})
\end{equation}
at every pole \(z_0\), and the exponent is positive for \(n\ge3\) and
\(p\ge4\).

By Corollary~\ref{cor:power-consequences}(c) the wave is untwisted and
\(\Gamma\) has rank one, and by Theorem~\ref{thm:A} its poles form one
orbit of \(\Gamma\). Place one of them at \(0\). Then
\(u(z)-(-1)^qu(-z)\) is invariant under \(\Gamma\), has no poles by
\eqref{eq:drift-parity}, is bounded at both ends and vanishes at \(0\);
so it is zero. Then all Laurent coefficients of odd index of \(u\) at
\(0\) vanish, in particular \(b_{p-1}=0\), which is false.
\end{proof}

In terms of \eqref{eq:pde} with zero integration constant: let
\(n\ge3\), let the order \(p+1\) be at least \(5\), and let
\(\varepsilon U_{xx}\) with \(\varepsilon\ne0\) be the only derivative
of even order. Then no travelling wave is meromorphic and nonconstant.
For example, travelling waves \(U=w(x-st)\) of
\[
 U_t+U^2U_x+\varepsilon U_{xx}+\delta U_{xxx}+U_{xxxxx}=0,
\]
with zero integration constant, satisfy \eqref{eq:main} with \(n=3\)
and \(P(\rho)=\rho^4+\delta\rho^2+\varepsilon\rho-s\), to which the
corollary applies when \(\varepsilon\ne0\).

\begin{remark}\label{rem:drift}
(a) In some cases the absence of waves has a local reason. For \(n=3\) and
\(P(D)=D^4+a_2D^2+a_1D+a_0\), the number \(F_J\) of
Section~\ref{sec:local} equals \(-a_1^2a_2/900\). So for
\(a_1a_2\ne0\) no Laurent series with a pole satisfies the equation,
and Proposition~\ref{prop:existence} alone shows that every meromorphic
solution is constant, as in \citep{KudryashovSinelshchikov2012SH} for
the Swift--Hohenberg equation without dispersion. For \(a_2=0\) the number \(F_J\) vanishes, and the
conclusion rests on Corollary~\ref{cor:drift}.

(b) The proof of Corollary~\ref{cor:drift} needs the exponent
\(p-1-q\) in \eqref{eq:drift-parity} to be positive; it is zero for
\(n=3\), \(p=2\) and negative for \(n=2\). For \(p=2\) the conclusion
does not hold: for \(n=3\), the function \(v=c\,G\), with \(G=k\,e^{kz}/(e^{kz}-1)\)
and \(c^2=-6\), solves \(v''-3kv'+2k^2v+v^3/3=0\), and for \(n=2\) the front of the
Korteweg--de Vries--Burgers equation is a wave of
\(v''+a_1v'+a_0v+v^2/2=0\) with \(a_1\ne0\) and
\(a_0=\pm6a_1^2/25\).
\end{remark}

For \(n=2\) and even \(p\ge4\) the conclusion of
Corollary~\ref{cor:drift} holds at least up to degree ten.

\begin{proposition}[\(n=2\), \(p\le10\)]\label{prop:drift-n2}
Let \(n=2\), let \(p\in\{4,6,8,10\}\), and let \(P(D)=P_e(D)+a_1D\),
where \(P_e\) is even and \(a_1\ne0\). Then \eqref{eq:main} has no
waves.
\end{proposition}

\begin{proof}
Suppose there is a wave; by Theorem~\ref{thm:A} it has one of the
forms (i)--(iii). Here \(m=1\) and \(q=p\). As in the proof of
Corollary~\ref{cor:drift}, the recurrence gives \(b_j=0\) for odd
\(j<p-1\) and
\(\Res_{z=0}u=b_{p-1}=pa_1/(p!-2K)\ne0\). So the wave is not elliptic.
For a rational wave \(u=\sum_{j=1}^{p}d_jz^{-j}\), the coefficients of
\(z^{-1}\) and \(z^{-2}\) in \eqref{eq:normalized} at infinity give
\(a_0d_1=0\) and \(a_0d_2-a_1d_1=Kd_1^2\), so
\(d_1\in\{0,-a_1/K\}\). But \(d_1=b_{p-1}\), and \(b_{p-1}\) is neither
of these numbers, because \((2-p)K<0<p!\). For the waves that are
rational in an exponential, an exact elimination in the space of
Theorem~\ref{thm:A}(ii), with all coefficients of \(P_e\) unknown,
shows that \(a_1=0\) (Appendix~\ref{app:computations}).
\end{proof}

We do not know a proof for all even \(p\ge4\).

\subsection{The Swift--Hohenberg equation with dispersion}\label{sec:sh}

Theorem~\ref{thm:A} turns the search for waves into a finite
computation. We carry it out for the equation
\[
 U_t+2U_{xx}-\sigma U_{xxx}+U_{xxxx}=\alpha U-U^3 ,
\]
the pure cubic case of the equation studied by
\citet{KudryashovSinelshchikov2012SH}. Travelling waves \(U=w(x-st)\)
satisfy \(P(D)w+w^3=0\) with
\[
 P(\rho)=\rho^4-\sigma\rho^3+2\rho^2-s\rho-\alpha ;
\]
this is \eqref{eq:main} for \(v=\sqrt3\,w\), so \(n=3\) and \(q=2\).
The operator is not even unless \(\sigma=s=0\). We determine all waves
that are rational in an exponential.

By the scaling of Section~\ref{sec:local}, a wave of \(P\) with rate
\(2\lambda\) comes from a wave with rate \(2\) of the operator
\(\widetilde P\) given by
\[
 \widetilde P(\rho)=\rho^4+a_3\rho^3+a_2\rho^2+a_1\rho+a_0,\qquad
 P(\rho)=\lambda^4\widetilde P(\rho/\lambda).
\]
We first find the waves with rate \(2\), with \(a_3,\dots,a_0\)
unknown, and then return to \(P\).

By Theorem~\ref{thm:A}(ii), a wave with rate \(2\) lies, after a
translation, in one of two spaces according to its character. We place
the pole at \(i\pi/2\). If the multiplier under \(z\mapsto z+i\pi\) is
\(-1\), the space is spanned by \(\sech z\) and \(D\sech z\); if the
multiplier is \(1\), it is spanned by \(1\), \(\tanh z\) and
\(D\tanh z=\sech^2z\). The normalized equation
is \(\widetilde P(D)u=120u^3\). We scale the amplitude so that the
leading coefficient at \(i\pi/2\) is \(i/2\) in the first space and
\(-1/2\) in the second. The squares of these numbers are
\(\mp\tfrac14\), so the coefficient \(120\) becomes \(\mp480\), and the
wave has one of the forms
\begin{equation}\label{eq:sh-spaces}
\begin{aligned}
 w_1&=(D-\xi)\,\tfrac12\sech z, &\qquad
 &\text{with}\quad\widetilde P(D)w_1=-480\,w_1^3,\\
 w_2&=\beta+\tfrac12\sech^2z-\tfrac{\xi}{2}\tanh z, &\qquad
 &\text{with}\quad\widetilde P(D)w_2=480\,w_2^3 .
\end{aligned}
\end{equation}
The parameter \(\beta\) is the midpoint of the two end values
\(\beta\mp\xi/2\) of \(w_2\) as \(\operatorname{Re}z\to\pm\infty\).

Both substitutions are finite computations. For \(w_2\), write
\(Y=(1+e^{2z})^{-1}\), so that \(\tanh z=1-2Y\), \(DY=-2Y(1-Y)\), and
\(w_2\) is a polynomial of degree two in \(Y\). Then
\(\widetilde P(D)w_2-480w_2^3\) is a polynomial of degree at most six
in \(Y\). The coefficient of \(Y^6\) vanishes by the normalization, the
coefficients of \(Y^5,\dots,Y^2\) determine \(a_3,\dots,a_0\) in turn
(Table~\ref{tab:sh}, top), and what remains is affine in \(Y\). An
affine function of \(Y\) is zero exactly when it vanishes at \(Y=0\)
and \(Y=1\), that is, at the two ends, where all derivatives of \(w_2\)
tend to \(0\). So the remaining conditions are
\(w_\pm(a_0-480\,w_\pm^2)=0\) for the two end values
\(w_\pm=\beta\mp\xi/2\). The end values are therefore equal, opposite,
or one of them is zero, which puts \((\beta,\xi)\) on one of the lines
\(\xi=0\), \(\beta=0\), \(\beta=\pm\xi/2\); on each line, inserting
\(a_0\) from the top of the table gives the equation in the lower part
of Table~\ref{tab:sh}. The four lines meet only at the origin
\(\beta=\xi=0\), which gives the wave \(\tfrac12\sech^2z\). For \(w_1\)
the same computation determines \(a_3,\dots,a_0\) and leaves one
polynomial condition on \(\xi\).

\begin{table}[ht]
\centering
\small
\begin{tabular}{@{}lll@{}}
\toprule
 & twisted, \(w_1\) & untwisted, \(w_2\)\\
\midrule
\(a_3\) & \(-14\xi\) & \(-14\xi\)\\
\(a_2\) & \(46\xi^2+10\) & \(46\xi^2-120\beta-20\)\\
\(a_1\) & \(-2\xi(7\xi^2+10)\) & \(-2\xi(7\xi^2-300\beta-20)\)\\
\(a_0\) & \(-14\xi^4-80\xi^2-11\) &
 \(1440\beta^2-120\beta\xi^2+480\beta-14\xi^4+160\xi^2+64\)\\
condition & \(\xi(\xi^2-1)(7\xi^2+17)=0\) & see below\\
\bottomrule
\end{tabular}

\medskip
\begin{tabular}{@{}ll@{}}
\toprule
line & remaining equation for \(w_2\) (necessary and sufficient)\\
\midrule
\(\beta=-\xi/2\) & \(\xi(\xi-4)(\xi^2-4)(7\xi-2)=0\)\\
\(\beta=\xi/2\) & \(\xi(\xi+4)(\xi^2-4)(7\xi+2)=0\)\\
\(\xi=0\) & \(\beta(30\beta^2+15\beta+2)=0\)\\
\(\beta=0\) & \(\xi(\xi^2-4)(7\xi^2+8)=0\)\\
\bottomrule
\end{tabular}
\caption{The forms \eqref{eq:sh-spaces} at rate \(k=2\). Top: the
coefficients of \(\widetilde P\) and the remaining condition. Bottom:
the remaining equation for \(w_2\) on each of the four lines.}
\label{tab:sh}
\end{table}

Altogether there are five values of \(\xi\) for the twisted form
\(w_1\) and fifteen pairs \((\beta,\xi)\) for the untwisted form
\(w_2\). To return to the physical equation, rescale by
\(\lambda\ne0\): the wave is
\(A\,w_j(\lambda z)\) for the operator
\(P(\rho)=\lambda^4\widetilde P(\rho/\lambda)\), with
\(\Gamma=(i\pi/\lambda)\ZZ\), that is, with rate \(k=2\lambda\), and with
\(A^2=480\lambda^4\) for \(w_1\) and \(A^2=-480\lambda^4\) for \(w_2\).
Comparing coefficients with \(P\) gives
\[
 \lambda^2a_2=2,\qquad \sigma=-\lambda a_3,\qquad
 s=-\lambda^3a_1,\qquad \alpha=-\lambda^4a_0 .
\]
Here \(a_2\ne0\) for all twenty parameter values, and
\(\sigma^2=2a_3^2/a_2\). Replacing \(z\) by \(-z\) turns \(w_1\) and
\(w_2\) into \(-w_1\) and \(w_2\) with \(-\xi\) in place of \(\xi\); so
the parameter values with \(\xi\) and \(-\xi\), taken with the factors
\(\lambda\) and \(-\lambda\), describe the same wave of the same
equation. As \(s\) and \(\alpha\) vary, the dispersion
coefficients that admit waves rational in an exponential are
\begin{equation}\label{eq:sh-sigma}
\begin{aligned}
 \text{twisted:}&\quad \sigma^2\in\Bigl\{0,\ 7,\ \frac{833}{89}\Bigr\},\\
  \text{untwisted:}&\quad \sigma^2\in\Bigl\{0,\ \frac{392}{71},\
 \frac{784}{127},\ \frac{1568}{239},\ \frac{392}{41},\
 \frac{392}{11}\Bigr\}.
\end{aligned}
\end{equation}
Each value of \(\sigma\) comes with the values of \(s\) and \(\alpha\)
determined by Table~\ref{tab:sh}. By Theorem~\ref{thm:A} there are no
further waves rational in an exponential, and by
Proposition~\ref{prop:not-even} an equation with \(\sigma\ne0\) has at
most one wave; each of the two values of \(\sigma\) with
\(\sigma^2=392/11\) occurs for two waves, which have different \(s\)
and \(\alpha\).

The list agrees with \citep{KudryashovSinelshchikov2012SH}, where the
equation contains an additional quadratic term \(bU^2\). The untwisted
waves with \(\sigma\ne0\) are its solutions with one pole per period,
at \(b=0\) in conditions (25) and (26) of that paper. The twisted waves
with \(\sigma\ne0\) are degenerations of its elliptic solutions with
poles of two types; those solutions satisfy the
condition\footnote{The condition is \(\alpha=2b^2/(9\kappa)\) for the
nonlinear terms \(\alpha U+bU^2+\kappa U^3\); here \(\kappa=-1\), and
\(\kappa=30\) in \citep{KudryashovSinelshchikov2012SH}.} under which
adding a constant to \(U\) removes the quadratic term and leaves no
constant term. The stationary waves are in Section~3 of that paper.

The list also illustrates the results of the preceding sections. By
Corollaries~\ref{cor:residue} and~\ref{cor:power-consequences}(b), every
elliptic wave and every twisted wave rational in an exponential has
\(\mathcal R(P)=0\), which for the Swift--Hohenberg operator reads
\[
 s=\frac{\sigma\,(392-39\sigma^2)}{1372};
\]
this is the speed of the elliptic solutions in
\citep{KudryashovSinelshchikov2012SH}. And the entries with
\(\sigma=0\) agree with Corollary~\ref{cor:drift}: for \(\sigma=0\) and
\(s\ne0\), the operator is an even polynomial minus \(sD\), so there
are no waves; accordingly, \(\sigma=-\lambda a_3=14\lambda\xi\) vanishes
only for \(\xi=0\), and then \(a_1=0\) and \(s=0\).

Few of these waves are real on the real axis. Suppose that a wave
\(w=A\,w_j(\lambda z)\) of the list is. At every pole \(w^2\sim-120\,(z-z_0)^{-4}\), so the leading coefficient is
purely imaginary. At a real pole it would be real; so \(w\) has no real
pole. Its poles come in conjugate pairs, with conjugate, hence
opposite, leading coefficients. An untwisted wave has the same leading
coefficient at all its poles; so no untwisted wave is real. For the
twisted waves with \(\xi^2=-17/7\) we have \(\lambda^2=-7/356\), so all
poles lie on one line parallel to the real axis; a conjugate pair of
poles would place this line on the real axis, where \(w\) has no pole;
so these waves are not real either. The twisted waves with \(\xi=0\) and \(\xi=\pm1\) are
real: there \(\lambda^2=1/5\) and \(1/28\), so \(\lambda\) and \(A\)
are real, and \(w_1\) is real on the real axis for real \(\xi\).

% ======================================================================
\section{Open questions}\label{sec:scope}

The bound of two waves is attained for \(p=4\) with \(n=2\) and \(n=3\)
(Sections~\ref{sec:two-waves} and~\ref{sec:n2}) and for \(p=6\) with
\(n=2\) (Section~\ref{sec:n2}). We do not know an equation with two
waves for \(n=3\) and \(p\ge6\), or for \(n=2\) and \(p\ge8\).

Theorems~\ref{thm:A} and~\ref{thm:B} describe all waves of a single
equation. The inverse question, which operators admit waves, is not
treated here beyond the examples of Sections~\ref{sec:example},
\ref{sec:power} and~\ref{sec:consequences}; for \(n=2\), the
operator--wave pairs with prescribed rate or lattice are counted in
\citep{MigitaBinomialCount2026}. How to select the real bounded waves
from these complex spaces, and how stable they are, remain open
questions.

\appendix

% ======================================================================
\section{Notation}\label{app:notation}

Table~\ref{tab:notation} lists the main symbols, with the place where
each is introduced.

\begin{table}[H]
\centering
\small
\begin{tabular}{@{}l@{\hspace{1em}}>{\raggedright\arraybackslash}p{0.56\textwidth}@{\hspace{1em}}l@{}}
\toprule
symbol & meaning & see\\
\midrule
\(D\); \(n\), \(m=n-1\) & \(d/dz\); power of the nonlinearity; phases \(\omega\in\mu_m\) & \eqref{eq:main}\\
\(p\), \(q=p/m\) & degree of \(P\); order of every pole & \eqref{eq:balance}\\
\((a)_k\) & rising factorial \(a(a+1)\cdots(a+k-1)\) & \eqref{eq:balance}\\
wave & nonconstant solution of \eqref{eq:main} that is meromorphic in \(\CC\) & Section~\ref{sec:intro}\\
\(\Wclass\) & functions that are rational, rational in an exponential, or elliptic & Section~\ref{sec:intro}\\
\(\Gamma\), \(\chi\) & group of twisted periods of a wave; its character & Definition~\ref{def:twisted}\\
twisted, untwisted & with nontrivial, trivial character & Section~\ref{sec:intro}\\
\(k\); ends; front & rate, \(\Gamma=(2\pi i/k)\ZZ\); \(\operatorname{Re}(kz)\to\pm\infty\); wave with different limits at the ends & Section~\ref{sec:intro}\\
\(a_{p-k}\) & coefficient of \(D^{p-k}\) in \(P\); it has weight \(k\) & Section~\ref{sec:local}\\
\(K\), \(c\) & \(K=(-1)^p(q)_p\) and \(c^m=-nK\) & \eqref{eq:normalized}\\
\(u\) & normalized wave, \(u(z)=v(z_0+z)/c=z^{-q}(1+O(z))\), \(P(D)u=Ku^n\) & \eqref{eq:normalized}\\
\(b_j\), \(I(j)\) & normalized Laurent coefficients; indicial polynomial & \eqref{eq:recurrence}\\
\(J\), \(h\) & for even \(p\): resonance index \(J=(n+1)q\); free coefficient \(h=b_J\) & Lemma~\ref{lem:resonance}, \eqref{eq:family}\\
\(f(z;h)\), \(F_J\) & series defined by the recurrence; \(P(D)f-Kf^n=-F_Jz^q\) & \eqref{eq:family}\\
\(\omega_i\), \(h_i\) & phase and free coefficient of a wave at its pole \(z_i\) & \eqref{eq:pole-germs}\\
\(S\), \(d(z)\) & pole set of a solution; distance to it & Section~\ref{sec:analytic}\\
\(g_r\), \(G=g_m\) & functions with simple poles on \(\Gamma=(2\pi i/k)\ZZ\) & \eqref{eq:periodic-form}\\
\(V_j^\chi\), \(g_\chi\) & functions of the second kind on a lattice, poles of order \(\le j\); element of \(V_1^\chi\) with residue \(1\) & Section~\ref{sec:spaces}\\
\(\wp(z;g_2,g_3)\), \(\wp(z;\Gamma)\) & Weierstrass function with \((\wp')^2=4\wp^3-g_2\wp-g_3\); that of the lattice \(\Gamma\) & Section~\ref{sec:two-waves}\\
\(\hatt\alpha\) & principal-part polynomial of a Laurent series \(\alpha\) & Section~\ref{sec:selectors}\\
\(E\), \(F\); \(\varphi\), \(\Phi\) & two expressions in \(u\); the combination \eqref{eq:phi-formal} on \(f(z;h)\) and on a wave & \eqref{eq:phi-formal}\\
\(H\), \(H_n\) & \(H(u)=uu''-\tfrac{q+1}{q}(u')^2\) and \(H_n=u^{n-2}H\) & \eqref{eq:seed}, Section~\ref{sec:quadratic}\\
\(\Delta_1\) & slope in \(h\) of the coefficient \([z^{p-2q-1}]\varphi\), for \(n\ge4\) & Proposition~\ref{prop:affine}\\
\([h^k]\), \([z^j]\) & coefficient of \(h^k\), of \(z^j\) in a series & Section~\ref{sec:selectors}\\
\(\psi\), \(\kappa\), \(\Delta_2\); \(X\) & leading coefficients of the terms in \(h^2\), for \(n=2,3\); \(\psi=\binom n2/(X-n)\) & \eqref{eq:psi}--\eqref{eq:Delta2}\\
\(a_*\) & leading coefficient of \(\hatt E\) for \(E=u\) & Proposition~\ref{prop:quadratic}\\
\(A_i\), \(B_i\) & \([z^i]\varphi=A_ih+B_i\) for the pair \((u,H_n)\) and \(i\le J-2\) & \eqref{eq:slopes}\\
\(T\); \(Q_P\) & entire function whose values at the poles do not contain \(h\); quadratic polynomial with \(Q_P(h)=0\) & Proposition~\ref{prop:quadratic}\\
\(c_0\), \(c_1\), \(c_{j-1}\) & numerical factors of \(A_0\), \(A_1\), \(A_{j-1}\) & Lemma~\ref{lem:slopes}\\
\(\ell_j\), \(\mathcal R(P)\) & coefficient of \(z^{j-p}\) in \(u^m\); the residue \(\mathcal R(P)=\ell_{p-1}\) & \eqref{eq:ell}, Corollary~\ref{cor:residue}\\
\(\Lambda_r\) & invariant function with \(\Lambda_r=z^{-r}+O(z)\) & Section~\ref{sec:power}\\
\bottomrule
\end{tabular}
\caption{Notation used throughout.}
\label{tab:notation}
\end{table}



% ======================================================================
\section{Elliptic functions of the second kind}\label{app:deferred}\label{app:second-kind}

We prove Lemma~\ref{lem:second-kind}. Recall that \(V_j^\chi\) is the
space of meromorphic functions \(w\) with
\(w(z+\gamma)=\chi(\gamma)w(z)\) for \(\gamma\) in the lattice
\(\Gamma\), whose poles lie in \(\Gamma\) and have order at most \(j\).

\begin{proof}[Proof of Lemma~\ref{lem:second-kind}]
A function \(w\) in \(V_0^\chi\) is entire, and \(|w|\) is periodic, so
\(w\) is bounded and hence constant; a nonzero constant has trivial
character. Thus \(V_0^\chi=0\) if \(\chi\) is not trivial.

For existence, let \(\sigma\) be the Weierstrass sigma function of
\(\Gamma\), and let \(\eta\colon\Gamma\to\CC\) be the quasi-period
homomorphism of the zeta function, \(\zeta(z+\gamma)=\zeta(z)+\eta(\gamma)\).
The function \(e^{az}\sigma(z-z_\chi)/\sigma(z)\) has the multiplier
\(\exp\bigl(a\gamma-\eta(\gamma)z_\chi\bigr)\) under \(z\mapsto z+\gamma\).
If \(\gamma_1,\gamma_2\) is a basis of \(\Gamma\), Legendre's relation
\(\eta(\gamma_1)\gamma_2-\eta(\gamma_2)\gamma_1=\pm2\pi i\) shows that
\((a,z_\chi)\mapsto\bigl(a\gamma_1-\eta(\gamma_1)z_\chi,\,a\gamma_2-\eta(\gamma_2)z_\chi\bigr)\)
is a linear bijection of \(\CC^2\); so \(a\) and \(z_\chi\) can be chosen
to realize any character. If \(\chi\) is not trivial, then
\(z_\chi\notin\Gamma\): otherwise the function would be a constant
multiple of some \(e^{bz}\), whose multipliers \(e^{b\gamma}\) have
finite order for all \(\gamma\) in a lattice only if \(b=0\), and then
\(\chi\) would be trivial. So the function has a simple pole at \(0\),
and dividing by its residue gives \(g_\chi\). The derivatives
\(D^ig_\chi\) have leading terms \((-1)^ii!\,z^{-i-1}\), so a suitable
combination of them has any prescribed principal part at \(0\). The
difference from a given element of \(V_j^\chi\) lies in \(V_0^\chi=0\).
This proves the basis property and the uniqueness.

For the trivial character, the residue at the only pole of an elliptic
function is zero, and the derivatives of \(\wp\) remove all higher
principal coefficients, leaving a constant.
\end{proof}

For \(m=2\) the functions \(g_\chi\) of the three characters of order
two are
\begin{equation}\label{eq:jacobi}
 g_{\chi_i}=\sqrt{\wp-e_i}\qquad(i=1,2,3),
\end{equation}
where \(e_1,e_2,e_3\) are the values of \(\wp\) at the three
half-periods of \(\Gamma\) and the root is chosen with leading term
\(z^{-1}\). Up to scale these are the Jacobi functions \(\cs\), \(\ds\),
\(\ns\); translation by the half-period that is called \(iK'\) in
Jacobi's notation turns them into multiples of \(\dn\), \(\cn\), \(\sn\)
\citep[\S\S22.4, 23.6]{DLMFonline}.

% ======================================================================
\section{The terms linear in the free coefficient}\label{app:slopes}

This appendix computes the constants of Lemma~\ref{lem:slopes}. Its
formulas, in particular \eqref{eq:ap-data} and \eqref{eq:first-odd},
are also checked by exact computation in the script
\mbox{\path{check_linear_terms.py}} (Appendix~\ref{app:computations}).
We keep the assumptions of the lemma: \(p\ge4\) is even, \(q\ge2\),
\(E=u\) and \(F=H_n\). Fix an operator \(P\) and a value of \(h\), and write
\[
 f=f(z;h),\qquad f_h=\frac{\partial f}{\partial h},\qquad
 \eta=H_n(f),\qquad \eta_h=\frac{\partial\eta}{\partial h};
\]
for series in \(z\), a prime denotes the derivative with respect to \(z\). By
Lemma~\ref{lem:degree} and the table \eqref{eq:leading-terms},
\(f_h=z^{nq}(1+O(z))\) and
\(\eta_h=J(J+1)z^{q-2}(1+O(z))\) are power series. The principal parts
of \(f\) and \(\eta\) do not depend on \(h\), so
\begin{equation}\label{eq:lin-phi}
 \frac{\partial\varphi}{\partial h}=\hatt E(D)\,\eta_h-\hatt F(D)\,f_h,
\end{equation}
where \(\hatt E\) and \(\hatt F\) are the principal-part polynomials of
\(f\) and \(\eta\). By \eqref{eq:slopes}, \(A_i\) is the coefficient of
\(z^i\) in \(\partial\varphi/\partial h\) for \(i\le J-2\), at every
value of \(h\).

\subsection{Residue formulas}

Let \(\alpha\) be a Laurent series with finite principal part
\(\sum_r\alpha_{-r}z^{-r}\), and \(g=\sum_kg_kz^k\) a power series. By
the definition of \(\hatt\alpha\),
\begin{align*}
 [z^0]\,\hatt\alpha(D)g&=\sum_r(-1)^{r-1}\alpha_{-r}\,g_{r-1}
 =\Res_{z=0}\alpha(z)\,g(-z),\\
 [z^1]\,\hatt\alpha(D)g&=\sum_r(-1)^{r-1}r\,\alpha_{-r}\,g_{r}
 =\Res_{z=0}\alpha'(z)\,g(-z).
\end{align*}
With \eqref{eq:lin-phi} this gives, for every value of \(h\),
\begin{equation}\label{eq:A-residues}
 A_0=\Res_{z=0}\bigl[f(z)\,\eta_h(-z)-\eta(z)\,f_h(-z)\bigr],\qquad
 A_1=\Res_{z=0}\bigl[f'(z)\,\eta_h(-z)-\eta'(z)\,f_h(-z)\bigr].
\end{equation}

\subsection{Even operators}

We first show that \(A_1\) contains no term \(a_{p-2}\).

\begin{lemma}\label{lem:even-A1}
If \(P\) is even, then \(A_1=0\).
\end{lemma}

\begin{proof}
For even \(P\) the recurrence preserves the parity of Laurent indices,
so \(b_j(h)=0\) for odd \(j\). Hence \(f_h(-z)=(-1)^{nq}f_h(z)\) and
\(\eta_h(-z)=(-1)^{q}\eta_h(z)\), and \((-1)^{nq}=(-1)^q\) because \(p\)
is even. So \eqref{eq:A-residues} gives
\[
 A_1=(-1)^q\,\Res_{z=0}\bigl[f'\eta_h-\eta'f_h\bigr].
\]

Let \(\mathcal L\) be the linearization of \(H_n\) at \(f\):
\[
 \mathcal L g=f^{\,n-1}g''-\frac{2(q+1)}{q}\,f^{\,n-2}f'\,g'+\theta\,g
\]
with a Laurent series \(\theta\). Then \(\eta_h=\mathcal Lf_h\), and
\(\eta'=\mathcal Lf'\) because \(H_n\) has constant coefficients. With
the Wronskian \(W=f'f_h'-f''f_h\) we get
\[
 f'\eta_h-\eta'f_h=f^{\,n-1}W'-\frac{2(q+1)}{q}\,f^{\,n-2}f'\,W .
\]
The residue of a derivative vanishes. Integrating by parts twice, and
using \((n-1)+2(q+1)/q=(J+2)/q\), we find
\[
 \Res\bigl[f'\eta_h-\eta'f_h\bigr]=-\frac{J+2}{q}\,\Res\bigl[f^{\,n-2}f'\,W\bigr]
 =\frac{J+2}{q}\,\Res\bigl[f_h\,\Theta\bigr],
\]
where
\[
 \Theta=(n-2)f^{\,n-3}(f')^3+3f^{\,n-2}f'f'' ;
\]
the last step holds because
\(f^{\,n-2}f'W=\bigl(f^{\,n-2}(f')^2f_h\bigr)'-f_h\Theta\).

Now \((f^{\,n})'''=n f^{\,n-1}f'''+n(n-1)\,\Theta\). Put
\(\mathcal M=P(D)-nKf^{\,n-1}\), and recall from
Section~\ref{sec:local} that \(P(D)f-Kf^n=-F_Jz^q\), where the number
\(F_J\) does not depend on \(h\). Differentiating this identity with
respect to \(h\) gives \(\mathcal Mf_h=0\), and differentiating it
three times with respect to \(z\) gives
\[
 \mathcal Mf'''=n(n-1)K\,\Theta-F_J\,(z^q)''' .
\]
The product of \(f_h\) with a polynomial has no residue, so
\(n(n-1)K\Res[f_h\Theta]=\Res[f_h\,\mathcal Mf''']\). Finally, for an even
operator \(P\), integration by parts gives
\(\Res[\alpha\,P(D)\beta]=\Res[\beta\,P(D)\alpha]\) for all Laurent
series \(\alpha,\beta\). Hence \(\Res[f_h\,\mathcal Mf''']
=\Res[f'''\,\mathcal Mf_h]=0\), and \(A_1=0\).
\end{proof}

By Lemma~\ref{lem:slopes}(a), \(A_1=c_1a_{p-1}^2+c'a_{p-2}\) with numbers
\(c_1,c'\). For the even operator \(D^p+a_{p-2}D^{p-2}\),
Lemma~\ref{lem:even-A1} gives \(c'a_{p-2}=0\), so \(c'=0\).

\subsection{An explicit family}

We now compute \(c_0\) and \(c_1\). For \(\varepsilon\ne0\) consider
the operator
\[
 P_\varepsilon(\rho)=\prod_{i=0}^{p-1}\bigl(\rho+(q+i)\varepsilon\bigr),
 \qquad a_{p-1}=\varepsilon\sum_{i=0}^{p-1}(q+i)
 =\tfrac12\,\varepsilon\,p\,(J-1).
\]
Put \(\zeta=e^{-\varepsilon z}\), so that
\(D=-\varepsilon\zeta\,d/d\zeta\). For a function \(g\) of \(\zeta\) we
have \(\bigl(D+(q+i)\varepsilon\bigr)\bigl[\zeta^{q+i}g\bigr]
=-\varepsilon\,\zeta^{q+i+1}\,dg/d\zeta\). Applying the factors with
\(i=0,\dots,p-1\) in turn gives
\[
 P_\varepsilon(D)\bigl[\zeta^{q}g(\zeta)\bigr]
 =\varepsilon^{p}\,\zeta^{q+p}\,\frac{d^{p}g}{d\zeta^{p}} .
\]
Hence, if \(Y(x)\) satisfies \(Y^{(p)}=KY^n\), then
\(u=(-\varepsilon)^q\zeta^q\,Y(\zeta-1)\) satisfies
\(P_\varepsilon(D)u=Ku^n\), because \(q+p=nq\) and \(p\) is even. The
operator \(D^p\) admits the wave \(x^{-q}\), so its series
\eqref{eq:family} is a formal solution of \(Y^{(p)}=KY^n\) for every
value \(\tilde h\) of the free coefficient, and by
Lemma~\ref{lem:degree} it equals
\(x^{-q}\bigl(1+\tilde hx^J+O(\tilde h^2)\bigr)\). Therefore, with
\(x=e^{-\varepsilon z}-1\),
\begin{equation}\label{eq:ap-family}
 (-\varepsilon)^q(1+x)^q\,x^{-q}\bigl(1+\tilde hx^J+O(\tilde h^2)\bigr)
\end{equation}
is a formal solution of \(P_\varepsilon(D)u=Ku^n\). It starts with
\(z^{-q}\), so it is the series \(f(z;h)\) of \(P_\varepsilon\), where
\(h\) is its coefficient of \(z^{J-q}\). The term linear in
\(\tilde h\) is
\((-\varepsilon)^q(1+x)^qx^{nq}\,\tilde h
=\varepsilon^Jz^{nq}(1+O(z))\,\tilde h\), because
\(x=-\varepsilon z(1+O(z))\) and \((-\varepsilon)^{q+nq}=\varepsilon^J\),
and the terms of higher degree in \(\tilde h\) are \(O(z^{2J-q})\). So
\(h=h_0+\varepsilon^J\tilde h\) for a constant \(h_0\). For \(\tilde h=0\) it is
the wave \(u=\bigl(\varepsilon/(e^{\varepsilon z}-1)\bigr)^q\) of
\(P_\varepsilon\).

We evaluate \eqref{eq:A-residues} at \(h=h_0\), that is, at
\(\tilde h=0\), where
\(\partial/\partial h=\varepsilon^{-J}\partial/\partial\tilde h\). By
\eqref{eq:seed-y}, \(H_n(f)=-q\,y^{-nq-1}y''\) with \(y=f^{-1/q}\), and
\eqref{eq:ap-family} gives
\(y=-x\,(1+\tilde hx^J)^{-1/q}/(\varepsilon(1+x))+O(\tilde h^2)\). Using
\(D=-\varepsilon(1+x)\,d/dx\), a direct computation yields, at
\(\tilde h=0\),
\begin{equation}\label{eq:ap-data}
\begin{aligned}
 f&=(-\varepsilon)^q(1+x)^q\,x^{-q}, \qquad
 \varepsilon^Jf_h=(-\varepsilon)^q(1+x)^q\,x^{nq},\qquad
 \eta=q\,\varepsilon^2(-\varepsilon)^{nq}(1+x)^{nq}\,x^{-nq-1},\\
 \varepsilon^J\eta_h&=\varepsilon^2(-\varepsilon)^{nq}(1+x)^{nq}\,x^{q-2}
 \bigl(J(J+1)+(2J^2+J+nq)\,x+J^2x^2\bigr).
\end{aligned}
\end{equation}
We compute the residues in the variable \(x\). Here
\(dz=-dx/(\varepsilon(1+x))\), and replacing \(z\) by \(-z\) replaces
\(x\) by \(-x/(1+x)\), so that a product \((1+x)^\mu x^\nu\)
becomes \((-1)^\nu(1+x)^{-\mu-\nu}x^\nu\). Thus
\[
 \varepsilon^Jf_h(-z)=(-1)^q(-\varepsilon)^q\,\frac{x^{nq}}{(1+x)^{J}},\qquad
 \varepsilon^J\eta_h(-z)=(-1)^q\varepsilon^2(-\varepsilon)^{nq}\,
 \frac{x^{q-2}\bigl(J(J+1)+qx-nqx^2\bigr)}{(1+x)^{J}},
\]
and \(f'=q\,\varepsilon(-\varepsilon)^q(1+x)^qx^{-q-1}\),
\(\eta'=q\,\varepsilon^3(-\varepsilon)^{nq}(1+x)^{nq}x^{-nq-2}(nq+1+x)\).
Since \((-\varepsilon)^{q+nq}=\varepsilon^J\), this gives
\begin{align*}
 A_0&=(-1)^{q+1}\varepsilon\,\Res_{x=0}
 \Bigl[\frac{J(J+1)+qx-nqx^2}{x^{2}(1+x)^{nq+1}}
 -\frac{q}{x\,(1+x)^{q+1}}\Bigr]\\
 &=(-1)^q\,\varepsilon\,J(J+1)(nq+1),\\
 A_1&=(-1)^{q+1}q\,\varepsilon^2\,\Res_{x=0}
 \Bigl[\frac{J(J+1)+qx-nqx^2}{x^{3}(1+x)^{nq+1}}
 -\frac{nq+1+x}{x^{2}(1+x)^{q+1}}\Bigr]\\
 &=(-1)^{q+1}\varepsilon^2\,\tfrac12\,q\,J(J+1)(nq+1)(nq+2).
\end{align*}
Both are nonzero. Dividing by \(a_{p-1}\) and by \(a_{p-1}^2\), and
using \(J/p=(n+1)/(n-1)\), we obtain the constants \(c_0\) and \(c_1\)
of Lemma~\ref{lem:slopes}(b).

\subsection{The first odd coefficient}

It remains to compute \(c_{j-1}\) and to show that it is not zero. Let
\(j\) be odd with \(3\le j\le p-1\), and let
\(P=D^p+a_{p-j}D^{p-j}\). The number \(c_{j-1}\) is the coefficient of
\(a_{p-j}\) in \(A_{j-1}\), so we compute modulo \(a_{p-j}^2\); we write
\(a=a_{p-j}\) for short. By Lemma~\ref{lem:degree}, modulo \(a^2\) and
\(h^2\),
\[
 f(z;h)=z^{-q}\bigl(1+\beta_j\,a\,z^{j}+h\,z^{J}+\gamma_j\,a\,h\,z^{J+j}\bigr)
\]
with numbers \(\beta_j\) and \(\gamma_j\). The recurrence
\eqref{eq:recurrence} at the indices \(j\) and \(J+j\) reads
\[
 I(j)\,\beta_j=(q)_{p-j},\qquad
 I(J+j)\,\gamma_j=-\frac{(nq)!}{(q+j)!}+(n-1)\,\frac{(nq)!}{q!}\,\beta_j .
\]
By the proof of Lemma~\ref{lem:resonance}, \(I(j)<0\) and \(I(J+j)>0\).
So \(\beta_j<0\) and \(\gamma_j<0\).

We use \eqref{eq:lin-phi} at \(h=0\). Let \(\Bil\) be the bilinear form
of \(H\), as in the proof of Proposition~\ref{prop:quadratic}. A direct
computation gives
\(2\Bil(z^{i-q},z^{nq})=\bigl((J-i)(J-i+1)-2ni\bigr)z^{p-2+i}\), and
Lemma~\ref{lem:seed} gives
\(2\Bil(z^{-q},z^{nq+i})=(J+i)(J+i+1)\,z^{p-2+i}\) and
\(H\bigl(z^{-q}(1+\beta_ja z^j)\bigr)=j(j+1)\beta_ja\,z^{j-2q-2}\).
Hence, at \(h=0\) and modulo \(a^2\),
\begin{gather*}
 \eta=j(j+1)\,\beta_ja\,z^{j-nq-2},\qquad
 \eta_h=J(J+1)\,z^{q-2}+\bigl(s_j\beta_j+(J+j)(J+j+1)\gamma_j\bigr)a\,z^{q-2+j},\\
 s_j=(n-2)\bigl(J(J+1)+j(j+1)\bigr)+(J-j)(J-j+1)-2nj ;
\end{gather*}
the three parts of \(s_j\) come from \(f^{\,n-2}\), from the factor
\(f^{\,n-3}f_hH(f)\), and from \(\Bil(f,f_h)\). The principal-part
polynomials are, modulo \(a^2\),
\[
 \hatt E(\rho)=\frac{(-\rho)^{q-1}}{(q-1)!}
 +\beta_ja\,\frac{(-\rho)^{q-j-1}}{(q-j-1)!},\qquad
 \hatt F(\rho)=j(j+1)\,\beta_ja\,\frac{(-\rho)^{nq+1-j}}{(nq+1-j)!},
\]
where the second term of \(\hatt E\) is absent if \(j\ge q\). Reading
off the coefficient of \(a\,z^{j-1}\) in \eqref{eq:lin-phi}, with
\(j\) odd and \(p\) even, we obtain
\begin{equation}\label{eq:first-odd}
\begin{aligned}
 (-1)^{q}c_{j-1}={}&-\binom{q+j-2}{j-1}\bigl(s_j\beta_j+(J+j)(J+j+1)\gamma_j\bigr)\\
 &+J(J+1)\binom{q-2}{j-1}\beta_j-j(j+1)\binom{nq}{j-1}\beta_j ,
\end{aligned}
\end{equation}
where \(\binom{q-2}{j-1}=0\) if \(j\ge q\). Up to this point \(n\ge2\)
is arbitrary.

Since \(\beta_j<0\) and \(\gamma_j<0\), we can write
\eqref{eq:first-odd} as
\begin{align*}
 (-1)^{q}c_{j-1}={}&|\beta_j|\Bigl[\binom{q+j-2}{j-1}s_j-J(J+1)\binom{q-2}{j-1}\Bigr]\\
 &+\binom{q+j-2}{j-1}(J+j)(J+j+1)\,|\gamma_j|
 +j(j+1)\binom{nq}{j-1}|\beta_j| .
\end{align*}
The last two terms are positive. So \((-1)^qc_{j-1}>0\) will follow
from
\begin{equation}\label{eq:first-odd-ineq}
 \binom{q+j-2}{j-1}\,s_j\ \ge\ J(J+1)\binom{q-2}{j-1},
\end{equation}
which we now prove for \(n=3\) and for \(n=2\). For \(n=3\) we have
\(J=4q\), \(j\le2q-1\) and
\[
 s_j=J(J+1)+j(j-5)+(J-j)(J-j+1)>J(J+1),
\]
because \(j(j-5)\ge-6\) and \(J-j\ge2q+1\ge5\). Since
\(\binom{q+j-2}{j-1}\ge\binom{q-2}{j-1}\), this proves
\eqref{eq:first-odd-ineq}. For \(n=2\) we have \(J=3q\),
\(3\le j\le q-1\) and
\[
 s_j=J(J+1)-j(6q+5-j)\ \ge\ J(J+1)\Bigl(1-\frac{2j}{3q}\Bigr)>0,
\]
because \(j(6q+5-j)\le j(6q+2)\) and \(J(J+1)=3q(3q+1)\). Moreover
\[
 \binom{q+j-2}{j-1}\Big/\binom{q-2}{j-1}
 =\prod_{i=0}^{j-2}\Bigl(1+\frac{j}{q-2-i}\Bigr)
 \ \ge\ 1+\frac{j(j-1)}{q},
\]
since each of the \(j-1\) factors is at least \(1+j/q\). The product
of the two lower bounds is
\[
 \Bigl(1+\frac{j(j-1)}{q}\Bigr)\Bigl(1-\frac{2j}{3q}\Bigr)
 =1+\frac jq\Bigl(j-1-\frac23-\frac{2j(j-1)}{3q}\Bigr)
 \ \ge\ 1+\frac jq\cdot\frac{j-3}{3}\ \ge\ 1,
\]
where we used \(j<q\) in the first inequality. This proves
\eqref{eq:first-odd-ineq} for \(n=2\). In both cases
\(c_{j-1}\ne0\), which completes the proof of
Lemma~\ref{lem:slopes}.

% ======================================================================
\section{Degrees one and two}\label{sec:low-order}

Theorems~\ref{thm:A} and~\ref{thm:B} assume \(p\ge3\). The one-pole
conclusion extends to the following cases of low degree.

\begin{proposition}\label{prop:low-order}
Every wave has exactly one pole modulo \(\Gamma\), and hence one of the
forms of Proposition~\ref{prop:spaces}, in each of the following cases:
\textup{(a)}~\(p=1\); \textup{(b)}~\(P(D)=D^2+a_0\);
\textup{(c)}~\(P(D)=D^2+a_1D+a_0\) and the wave is elliptic. In case
\textup{(c)} necessarily \(a_1=0\).
\end{proposition}

\begin{proof}
(a) Since \(p\) is odd, the normalized series is unique; apply
Corollary~\ref{cor:one-germ}.

(b) Pole balance leaves \((n,q)=(2,2)\) or \((3,1)\), and \(J=2q+2\).
The energy
\[
 \mathcal E=\frac{(v')^2}{2}+\frac{a_0v^2}{2}+\frac{v^{n+1}}{n(n+1)}
\]
is constant along every solution. At a pole write
\(v=cz^{-q}\bigl(1+\sum_{j\ge1}b_jz^j\bigr)\) with \(h=b_J\). The
constant \(\mathcal E\) equals the constant term of its Laurent
expansion at the pole, which gives
\[
 \frac{\mathcal E}{c^2}=-q(2q+3)\,h+\mathcal E_0,
\]
where \(\mathcal E_0\) depends only on \(b_1,\dots,b_{J-1}\), hence
only on \(P\). Indeed, the coefficient of \(h\) is \(-q(q+2)\) from
\((v')^2/2\) plus \(-q(q+1)\) from the last term, while the \(h\)-term
of \(a_0v^2/2\) has a positive power of \(z\). In both admissible cases
\(c^2\) is the same for all leading coefficients. Hence \(h\) is the
same at all poles, and Corollary~\ref{cor:one-germ} applies.

(c) Pole balance again leaves \((n,q)=(2,2)\) or \((3,1)\). At each
pole, the recurrence gives \(c^m=-nq(q+1)\) and \(b_1=-a_1/(n+3)\).
The elliptic function \(v^m\) has a double pole at each pole of \(v\),
with residue
\[
 \Res v^m=mc^mb_1=\frac{mnq(q+1)}{n+3}\,a_1,
\]
the same at every pole. The residues in a period cell add up to zero,
so \(a_1=0\), and (b) applies.
\end{proof}

Part (b) is the case \(p=2\) of Lemma~\ref{lem:energy} below, since
\(I'(J)=2q+3\) when \(p=2\). In two respects degree two differs from
higher degree.

\emph{Infinitely many waves.}\label{sec:order-two} A single equation
can have infinitely many waves, so the bound of Theorem~\ref{thm:B}
requires \(p\ge3\). For \(n=2\) and \(P(D)=D^2+a_0\), the functions
\[
 v=-12\,\wp\bigl(z;\tfrac{a_0^2}{12},g_3\bigr)-a_0,\qquad
 g_3\in\CC,\quad 27g_3^2\ne\bigl(\tfrac{a_0^2}{12}\bigr)^3,
\]
solve \(v''+a_0v+v^2/2=0\), by the identity \(\wp''=6\wp^2-g_2/2\).
Their free coefficient is \(h=g_3/28\), so different values of \(g_3\)
give different normalized series, and the corresponding waves are not
related by translation or phase change. They still have one pole modulo
\(\Gamma\), as Proposition~\ref{prop:low-order}(b) asserts.

\emph{Waves outside \(\Wclass\).} Equations with \(a_1\ne0\) can have
meromorphic solutions outside \(\Wclass\)
\citep[Lemmas~4.4 and~4.5]{ConteNgWu2015}. For \(k\ne0\) and all
constants \(\zeta_0\) and \(g_3\), the function
\(v=-12k^2e^{-2kz}\wp(e^{-kz}-\zeta_0;0,g_3)\) is a meromorphic
solution of \(v''+5kv'+6k^2v+v^2/2=0\)
\citep{AblowitzZeppetella1979}, \citep[Lemma~4.4]{ConteNgWu2015}. For
\(g_3\ne0\) it lies outside \(\Wclass\), because it has too many poles:
they are the points with \(e^{-kz}\in\zeta_0+\Lambda\), where
\(\Lambda\) is the period lattice of \(\wp(\cdot\,;0,g_3)\), and their
number in \(|z|\le r\) grows exponentially with \(r\), whereas a
function of \(\Wclass\) has \(O(r^2)\) poles in \(|z|\le r\). The
expression
\((v'+2kv)^2+v^3/3=-144\,g_3k^6e^{-6kz}\) has no poles, and for
\(g_3\ne0\) it is not constant. So the hypothesis on the values at
the poles in Lemma~\ref{lem:finite-values} cannot be dropped.

% ======================================================================
\section{Polynomial nonlinearities}\label{app:general}

Moving a constant solution of a general polynomial equation to \(0\)
and absorbing the linear term into \(P\) gives
\begin{equation}\label{eq:B}
 P(D)v+B(v)=0,\qquad
 B(v)=\frac{v^n}{n}+B_*(v),
\end{equation}
where \(B_*\) is a polynomial that contains only the powers
\(v^2,\dots,v^{n-1}\). The leading balance at a pole is unchanged:
\(p=mq\) and \(c^m=-nK\). The Laurent recurrence has the same indicial
polynomial \(I\); the lower terms of \(B\) change only the right-hand
sides, which now also depend on the choice of \(c\). The phase changes
that preserve \eqref{eq:B} form the group \(\mu_{m'}\), where \(m'\) is
the largest divisor of \(m\) such that \(B\) contains only powers
\(v^j\) with \(j\equiv1\pmod{m'}\). For the pure power, \(m'=m\).

For an even operator \(P(D)=D^p+\sum_{0\le i<p/2}a_{2i}D^{2i}\), put
\[
 \mathcal Q_i(v)=\sum_{a=0}^{i-2}(-1)^a\,v^{(a+1)}v^{(2i-1-a)}
 +\frac{(-1)^{i-1}}{2}\bigl(v^{(i)}\bigr)^2 ,
\]
so that \(D\mathcal Q_i(v)=v'v^{(2i)}\); for \(i=1\) the sum is empty.
The energy
\begin{equation}\label{eq:energy}
 \mathcal E(v)=\mathcal Q_{p/2}(v)+\sum_{1\le i<p/2}a_{2i}\,\mathcal Q_i(v)
 +\frac{a_0}{2}v^2+\int_0^vB(x)\,dx
\end{equation}
satisfies \(D\mathcal E(v)=v'\bigl(P(D)v+B(v)\bigr)\), so it is constant
along every solution. This first integral fixes the free coefficient,
as in \citet*[Lemma~4]{EremenkoLiaoNg2009}. \citet{Demina2019} classifies the meromorphic integrals of autonomous
equations with two dominant monomials.

\begin{lemma}[The energy fixes the free coefficient]\label{lem:energy}
Let \(P\) be even, and fix a root \(c\) of \(c^m=-nK\). The Laurent
recurrence with leading coefficient \(c\) is compatible at \(J=p+2q\),
and its formal solutions form a one-parameter family
\(v=cz^{-q}\bigl(1+\sum_{j\ge1}b_jz^j\bigr)\) with free coefficient
\(h=b_J\). Along this family
\begin{equation}\label{eq:energy-h}
 \mathcal E(v)=-qc^2I'(J)\,h+\mathcal E_0,\qquad
 I'(J)=(q+1)_p\sum_{i=1}^{p}\frac{1}{q+i}>0,
\end{equation}
where \(\mathcal E_0\) depends only on \(P\), \(B\) and \(c\). Hence a
leading coefficient and a value of the energy determine the Laurent
series.
\end{lemma}

\begin{proof}
Construct the series up to index \(J-1\). For this truncated series the
lowest term of \(P(D)v+B(v)\) that can be nonzero is \(\Omega z^q\), and
the residue of \(v'(P(D)v+B(v))\) is then \(-qc\,\Omega\). This is the
residue of the derivative \(D\mathcal E\), so it vanishes, and
\(\Omega=0\). The coefficient \(b_J\) is therefore free, and all later
coefficients are determined.

Every term of \(\mathcal E\) has a pole of order at most \(J\). A
factor \(h\) comes with the power \(z^J\). So \(h\) enters the constant
term only through terms of pole order \(J\) in which all other factors
are at leading order, and it enters linearly. To compute the coefficient of \(h\), perturb
\(v=cz^{-q}(1+\varepsilon z^j)\). The top
linear term of the equation is \(\varepsilon cI(j)z^{j-q-p}\); write the
corresponding term of \(\mathcal E\) as \(\varepsilon L(j)z^{j-J}\). The
identity \(D\mathcal E=v'(P(D)v+B(v))\) gives \((j-J)L(j)=-qc^2I(j)\)
as polynomials in \(j\), so \(L(J)=-qc^2I'(J)\). The formula for
\(I'(J)\) follows by differentiating the product in \(I(j)\), whose
factors at \(j=J\) are \(q+p,\dots,q+1\). Finally, since
\(D\mathcal E=0\) formally, the constant Laurent coefficient of
\(\mathcal E\) is its value.
\end{proof}

The theorem below needs Lemma~\ref{lem:germs} for the solutions of
\eqref{eq:B} and for subgroups of \(\mu_m\). The proof of that lemma
uses only Lemma~\ref{lem:distance} and the finiteness of \(\mu_m\). So
the lemma holds for every nonconstant meromorphic solution \(v\) of
\eqref{eq:general-eq} and every finite group \(U\subset\CC^*\), with
\(\mu_m\) replaced by \(U\) and \(\Gamma\) by the group
\[
 \Gamma_U=\{\gamma\in\CC:\ v(z+\gamma)=\omega v(z)
 \text{ for some }\omega\in U\};
\]
since \(v\not\equiv0\), the multiplier \(\omega=\chi(\gamma)\) is
unique, and \(\chi\colon\Gamma_U\to U\) is a homomorphism.

\begin{theorem}[Polynomial nonlinearity; cf.\ \citep{Eremenko2006,ConteNgWong2012,Wu2014}]\label{thm:general}
Let \(B\) have the symmetry group \(\mu_{m'}\), and assume that \(p\) is
odd or \(P\) is even. Then every nonconstant meromorphic solution of
\eqref{eq:B} belongs to \(\Wclass\). Put \(\mathcal H=\mu_{m'}\) if
\(p\) is odd, and \(\mathcal H=\mu_{m'}\cap\{\pm1\}\) if \(P\) is even.
Two poles whose leading coefficients have ratio \(\omega\in\mathcal H\)
differ by a twisted period with multiplier \(\omega\), and the poles
form at most \(m/|\mathcal H|\) cosets of \(\Gamma_{\mathcal H}\).
\end{theorem}

\begin{proof}
For odd \(p\) the recurrence determines the Laurent series from its
leading coefficient. For even \(P\) the energy of a solution is
constant, so by Lemma~\ref{lem:energy} each leading coefficient carries
at most one Laurent series among the poles of the solution.

Let \(z_1,z_2\) be poles with \(c_2=\omega c_1\), \(\omega\in\mathcal
H\), and compare \(v(z_1+z)\) with \(\omega^{-1}v(z_2+z)\). The second
function is again a solution, because \(\omega\) is a symmetry, and the
two have the same leading coefficient. For odd \(p\) this identifies
the two Laurent series. For even \(P\), the symmetry gives
\(\mathcal E(\omega v)=\omega^2\mathcal E(v)\), since every quadratic
term scales by \(\omega^2\) and
\(\int_0^{\omega v}B=\omega^2\int_0^vB\); for \(\omega=\pm1\) the two
energies agree, and the energy lemma identifies the series. By the
identity theorem, \(z_2-z_1\) is a twisted period with multiplier
\(\omega\).

The \(m\) leading coefficients form \(m/|\mathcal H|\) orbits under
\(\mathcal H\), and two poles whose leading coefficients lie in the same
orbit differ by an element of \(\Gamma_{\mathcal H}\), as just shown.
Lemma~\ref{lem:germs}, in the form stated before the theorem with
\(U=\mathcal H\), gives membership in \(\Wclass\) and the bound on the
number of cosets.
\end{proof}

For the pure power and an even operator with \(n\le3\) we have
\(\mu_m\subset\{\pm1\}\); so the theorem gives one pole orbit, as
Lemma~\ref{lem:two-germs} does.

\begin{remark}
For an even operator and a polynomial nonlinearity, the reflection
argument of Proposition~\ref{prop:reflection} applies whenever
\((-1)^q\in\mu_{m'}\): the reflected series has the same leading
coefficient and the same free coefficient, because \(J\) is even. So
reflection symmetry holds for every even \(q\), and for odd \(q\) when
\(B\) is odd. This alone does not relate different orbits of leading
coefficients.
\end{remark}

% ======================================================================
\section{Computer checks}\label{app:computations}

The theorems are proved in the text for all \(n\) and \(q\). Exact
computations with Python~3 and SymPy check the displayed identities,
the constants of Sections~\ref{sec:selectors} and~\ref{sec:proof} and
of Appendix~\ref{app:slopes}, and the examples. The numbers quoted in
the examples without derivation, such as the entries of
Table~\ref{tab:sh}, the number \(F_J\) in Remark~\ref{rem:drift}(a) and
the coefficient \(A_5\) in Section~\ref{sec:not-even}, are obtained by
finite computations that can be done by hand. The only statement of the
paper that rests on a computer elimination is
Proposition~\ref{prop:drift-n2}; the script
\mbox{\path{check_drift_n2.py}} carries out the elimination for
\(p=4,6,8,10\) and shows that the statement does not hold for \(p=2\). The
scripts are the ancillary files (\path{anc/}); the command
\texttt{python3 run\_all.py} in that folder runs all of them, and
\path{anc/README.md} describes each script.

The official repository for the manuscript and ancillary materials is
\url{https://github.com/migita/meromorphic-waves-structure}.

\section*{Use of AI tools}

{\small
This paper was prepared with substantial use of generative AI, disclosed
as recommended by the Leiden Declaration on Artificial Intelligence and
Mathematics.\footnote{\url{https://leidendeclaration.ai/}, 2 June 2026.}
The main mathematical ideas were produced by \texttt{gpt-6-astra}
(OpenAI) in dialogue with the author; parts of the proofs, of the
verification and of the preparation of the text were carried out with
Claude Fable~5.1 and Claude Opus~5.5 (Anthropic). No automated system
is an author; responsibility for the results and for the attribution of
prior work rests with the author.\par}

\begingroup
\raggedright
\bibliographystyle{abbrvnat}
\bibliography{references}
\endgroup

\end{document}
